\documentclass[10pt,a4paper]{article}

% ----------------- Paquetes -------------------%
\usepackage[T1]{fontenc}
\usepackage[utf8]{inputenc}
\usepackage[a4paper,margin=2.5cm]{geometry}
\usepackage{mathptmx}             
\usepackage{changepage} 
\usepackage{setspace}
\usepackage[spanish]{babel}
\usepackage[labelfont=bf,labelsep=space]{caption}
\usepackage{amssymb}
\usepackage{amsmath}
\usepackage{ebgaramond}
\usepackage{placeins}
\usepackage{etoolbox}
\usepackage{setspace}
\usepackage{parskip} 
\usepackage{tcolorbox}
\usepackage{needspace}
\usepackage{xparse}
\usepackage{ocgx2}
\usepackage{hyperref}
\usepackage{hypcap}
\usepackage{soul}\sethlcolor{yellow}% resaltado amarillo: \hl{texto} (Panel TCEE, ⌃H)



%%%%%%%%%%%%%%%%%%%%%%%%%%%%%%%%%%%%%%%%%%%%%%%%%%%%%%%%%%%%%%%%
% --- Fija primero tus valores globales "buenos" ---
\setstretch{1.2}                 % Interlineado global
\setlength{\parskip}{0.7\baselineskip}
\setlength{\parindent}{0pt}
\newlength{\GlobalParskip}
\setlength{\GlobalParskip}{\parskip}

\newcounter{eqnum}[subsection]
\renewcommand{\theeqnum}{\arabic{eqnum}}
\newcounter{alphaitem}
\newcounter{numitem}

%% ------------------- Recuadros----------------------%%
\newcommand{\puntosclave}[1]{%
  \begin{tcolorbox}[
    colframe=black,
    title={Puntos clave},
    before upper={\setlength{\parindent}{0.5cm}\setlength{\parskip}{0.5\baselineskip}}
  ]
  #1
  \end{tcolorbox}
}

\newcommand{\modificaciones}[1]{
  \begin{tcolorbox}[
    colback=white,
    colframe=red,
    title={Modificaciones a realizar},
    before upper={\setlength{\parindent}{0.5cm}\setlength{\parskip}{0.5\baselineskip}}
  ]
  #1
  \end{tcolorbox}
}

%% ------------------- Preguntas TEST----------------------%%
% Opciones: radio (single-choice)
\NewDocumentCommand{\OptRadio}{m m m}{%
  \noindent\CheckBox[radio,name=#1,value=#2,width=1.2ex,height=1.2ex]{}%
  \;\textbf{#2.}~#3\par
}

% Opciones: checkbox (multi-choice)
\NewDocumentCommand{\OptCheck}{m m}{%
  \noindent\CheckBox[name=#1,width=1.2ex,height=1.2ex,bordercolor={0 0 0}]{}%
  \;#2\par
}

% Pregunta single-choice (radio)
% Args: 1=id, 2=enunciado, 3=opciones, 4=solución+justificación, 5=imagen (o {}), 6=origen
\NewDocumentCommand{\PreguntaTestSingle}{m m m m m m}{%
  \par\medskip
  \noindent\textbf{#1.}~#2\par\smallskip

    % Imagen (si no es {})
    \IfBlankTF{#5}{}{%
      \begin{center}
        \includegraphics[width=0.75\linewidth]{#5}
      \end{center}
    }

    \begin{Form}
      #3
    \end{Form}

    \par\smallskip
    \switchocg{sol-#1}{\fbox{Mostrar / ocultar solución}}%
    \hfill{\footnotesize #6}\par\smallskip

    \begin{ocg}{Solución}{sol-#1}{0}
      \noindent #4
    \end{ocg}
  \par\medskip
}

% Pregunta multi-choice (checkboxes)
\NewDocumentCommand{\PreguntaTestMulti}{m m m m m m}{%
  \par\medskip
  \noindent\textbf{#1.}~#2\par\smallskip

    % Imagen (si no es {})
    \IfBlankTF{#5}{}{%
      \begin{center}
        \includegraphics[width=0.75\linewidth]{#5}
      \end{center}
    }
    \begin{Form}
      #3
    \end{Form}

    \par\smallskip
    \switchocg{sol-#1}{\fbox{Mostrar / ocultar solución}}%
    \hfill{\footnotesize #6}\par\smallskip

    \begin{ocg}{Solución}{sol-#1}{0}
      \noindent #4
    \end{ocg}
  \par\medskip
}


%% ------------------- CITA ----------------------%%
\NewDocumentEnvironment{cita}{ O{} O{} +b }{%
  \begin{adjustwidth}{2cm}{0pt}
  \raggedright
  #3\par
  \raggedleft
  \if\relax\detokenize{#1#2}\relax
    % nada
  \else
    #1\if\relax\detokenize{#1}\relax\else\if\relax\detokenize{#2}\relax\else, \fi\fi#2
  \fi
  \end{adjustwidth}
}{}



%% ------------------- Listas----------------------%%
    % --- Lista alfabética ---
    \newenvironment{la}
      {\begin{list}{\alph{alphaitem})}{%
          \usecounter{alphaitem}%
          \setlength{\leftmargin}{1cm}%
          \setlength{\parskip}{3pt}% 
          \setlength{\itemsep}{0pt}% ← espacio entre ítems
          \setlength{\parsep}{0pt}%  ← párrafos dentro de un ítem
      }}
      {\end{list}}
      
    % --- Lista numérica ---
    \newenvironment{lnum}
      {\begin{list}{\arabic{numitem}.}{%
          \usecounter{numitem}%
          \setlength{\leftmargin}{1cm}%
          \setlength{\parskip}{3pt}% 
          \setlength{\itemsep}{0pt}% ← espacio entre ítems
          \setlength{\parsep}{0pt}%  ← párrafos dentro de un ítem
      }}
      {\end{list}}

%% ------------------- Preguntas Test ----------------------%%
      \usepackage{xparse} % si ya lo tienes, no lo repitas

    \NewDocumentEnvironment{test}{m m o}{%
      \begin{tcolorbox}[
        colback=white,
        colframe=black!15,
        boxrule=0.6pt,
        arc=2mm,
        left=2mm,right=2mm,top=2mm,bottom=2mm
      ]
      \centering
      \includegraphics[width=\linewidth]{#1}
      \end{tcolorbox}
    
      \vspace{2mm}
    
      % Caja SOLO para texto (SÍ partible y mantiene márgenes al partir)
      \begin{tcolorbox}[
        enhanced,
        breakable,
        colback=black!3,
        colframe=black!25,
        boxrule=0.6pt,
        arc=2mm,
        left=2mm,right=2mm,top=1.5mm,bottom=1.5mm
      ]
      \textbf{Respuesta:} \textbf{#2}%
      \IfValueT{#3}{\par\smallskip \textbf{Justificación:} #3}
      \end{tcolorbox}
    }{}
      
% ------------------- Imágenes----------------------%
    \graphicspath{{Img/}}% Rutas por defecto 
    % --- Imagen adaptativa ---
    % Registros de dimensión definidos una sola vez
    \newdimen\imgcap
    \newdimen\imgmin
    \newdimen\imgmax
    \newdimen\imgremain
    \newdimen\imgh
    \newdimen\imgneed
    
    \newcommand{\imagenfit}[3]{%
      \addvspace{10pt}% separación antes
      \par\begingroup
        % Parámetros de composición (asignaciones, no \newdimen)
        \imgcap=1\baselineskip            % alto estimado de título+fuente
        \imgmin=0.15\textheight
        \imgmax=0.15\textheight
        \imgremain=\pagegoal
        \ifdim\pagegoal=\maxdimen \imgremain=0pt\fi
        \advance\imgremain by -\pagetotal
        \advance\imgremain by -\imgcap
        % Decisión de altura destino
        \ifdim\imgremain<\imgmin
          \newpage
          \imgh=0.20\textheight
        \else
          \imgh=\imgremain
          \ifdim\imgh>\imgmax \imgh=\imgmax \fi
        \fi
        % Reservar espacio para evitar cortes del bloque completo
        \imgneed=\imgh \advance\imgneed by \imgcap
        \Needspace{\imgneed}
        % Cuerpo
        \noindent\begin{minipage}{\textwidth}
          \centering
          \captionof{figure}{\textemdash\ #2}\par\vspace{0.3em}% título arriba
          \includegraphics[width=\linewidth,height=\imgh,keepaspectratio]{\detokenize{#1}}\par
          \vspace{0.3em}\small\textit{Fuente: }#3% fuente abajo
        \end{minipage}%
      \endgroup
      \par\addvspace{10pt}%
    }

%------------ Bloque de RECUADRO ESPECIAL -------------%
    \newenvironment{specialblock}[1]{%
      \begin{quote} % crea sangría a izquierda y derecha
      \small % texto más pequeño
      \noindent\textbf{#1}\\[-0.3em] % título en negrita
      \rule{\linewidth}{0.4pt}\\[0.6em]}{
      \end{quote}}
      
%-------------------------- Portada -----------------------%
\newcommand{\oldtitle}[1]{\def\OldTitle{#1}}
\newcommand{\changerationale}[1]{\def\ChangeWhy{#1}}

\newcommand{\changebox}{%
\begin{tcolorbox}[title={Historial de cambios del tema}]
\textbf{Título anterior:} \OldTitle\par
\textbf{Motivo del cambio:} \ChangeWhy
\end{tcolorbox}
}

%-------------------------- Author Font -----------------------%
    \newcommand{\authorfont}[1]{{\fontfamily{EBGaramond-TLF}\selectfont #1}}

%-------------------------- Anexos -----------------------%
    \makeatletter
    \newcounter{anexo}
    \renewcommand{\theanexo}{\Roman{anexo}}
    \newcommand{\anexo}[1]{%
      \clearpage
      \phantomsection
      \refstepcounter{anexo}%
      \noindent\textbf{Anexo \theanexo.- }#1\par\medskip
      \addcontentsline{stc}{section}{Anexo \theanexo.- #1}%
    }
    \makeatother

    %-------------------------- Lista de figuras (personal) -----------------------%
\makeatletter
\newcommand{\MyListOfFigures}{%
  \clearpage
  \phantomsection
  {\centering\bfseries\Large Índice de figuras\par}%
  \medskip
  % Entrada en nuestro índice de contenido (stc)
  \addcontentsline{stc}{section}{Índice de figuras}%
  % Recuperación del listado (sin cabecera propia)
  \begingroup
    \parindent=0pt\parskip=2pt
    \@starttoc{lof}%
  \endgroup
}
\makeatother

%-------------------------- Bibliografía -----------------------%
\makeatletter
% Contador para las referencias
\newcounter{myref}
% Cabecera de la bibliografía, entra en el índice 'stc'
\newcommand{\MyBibliography}{%
  \clearpage
  \phantomsection
  {\centering\bfseries\Large Bibliografía y referencias\par}%
  \medskip
  \addcontentsline{stc}{section}{Bibliografía y referencias}%
  \setcounter{myref}{0}%
}
\newcommand{\MyBibItem}[4]{%
  \refstepcounter{myref}%
  \noindent[\themyref]\ #1.\ \emph{#2}. #3. #4\par
}
\makeatother

%--------- Establecer cuatro niveles de índice --------%
    \setcounter{tocdepth}{4}   
    \setcounter{secnumdepth}{4}% numera hasta \paragraph

%%%%%%%%%%%%%%%%%%%%%%%%%%%%%%%%

\makeatletter

    %--------- Interlineado Global --------%
    \edef\GlobalBaseLineStretch{\baselinestretch}
    
    %--------- Espaciado de seccinones --------%
    \renewcommand\section{\@startsection{section}
      {1}{\z@}
      {2.5ex \@plus 1ex \@minus .2ex} % espacio antes
      {1.5ex \@plus .2ex}             % espacio después
      {\normalfont\Large\bfseries}    % formato del título
    }
    \renewcommand\subsection{\@startsection{subsection}{2}{\z@}
      {3ex \@plus 0.8ex \@minus .2ex}
      {1.5ex \@plus .2ex}
      {\normalfont\large\bfseries}
    }
    \renewcommand\subsubsection{\@startsection{subsubsection}{3}{\z@}
      {3ex \@plus 0.5ex \@minus .2ex}
      {1ex \@plus .2ex}
      {\normalfont\normalsize\bfseries}
    }
    \renewcommand\paragraph{\@startsection{paragraph}{4}{\z@}%
    {-2ex \@plus -0.5ex \@minus -.2ex}% espacio antes
    {0.8ex \@plus .2ex}% espacio después
    {\normalfont\normalsize\itshape}}% ← cursiva en lugar de negrita

    %--------- Ecuaciones --------%   
    % Variables globales para almacenar títulos y notas
    \gdef\leq@current@section{}%
    \gdef\leq@current@subsection{}%
    \gdef\leq@last@printed@section{}%
    \gdef\leq@last@printed@subsection{}%
    
    % Comando para añadir nota (invisible en el cuerpo)
    \newcommand{\eqnote}[1]{%
      \addcontentsline{leq}{leqnote}{#1}%
    }
    
    % Comando para ecuaciones
    \newcommand{\eqblock}[2]{%
      % Verificar si necesitamos imprimir el título de sección
      \ifx\leq@current@section\leq@last@printed@section\else
        \addcontentsline{leq}{leqsection}{\leq@current@section}%
        \global\let\leq@last@printed@section\leq@current@section
        \gdef\leq@last@printed@subsection{}% Reset subsección al cambiar sección
      \fi
      % Verificar si necesitamos imprimir el título de subsección
      \ifx\leq@current@subsection\leq@last@printed@subsection\else
        \ifx\leq@current@subsection\@empty\else
          \addcontentsline{leq}{leqsubsection}{\leq@current@subsection}%
          \global\let\leq@last@printed@subsection\leq@current@subsection
        \fi
      \fi
      % Ecuación
      \refstepcounter{eqnum}%
      \begin{equation*}
        #1\tag{E.\theeqnum}
      \end{equation*}%
      \addcontentsline{leq}{leqt}{[E.\theeqnum] -- #2}%
      \addcontentsline{leq}{leqm}{#1}%
    }
    
    %%% ----- Índice de ecuaciones ----------%%%
    % Tamaño para []
    \newcommand{\leqIndexFont}{\normalsize}
    \newcommand{\leqTitleFont}{\tiny}
    \newcommand{\leqNoteFont}{\tiny}
    % Cabecera de sección 
    \newcommand*\l@leqsection[2]{%
      \addvspace{25pt}% Mayor separación antes de nueva sección
      \noindent\leqIndexFont
      \makebox[\linewidth]{\rule{.38\linewidth}{0.4pt}\ \ \textbf{  \MakeUppercase{#1}}\ \ \rule{.38\linewidth}{0.4pt}}%
      \par\vspace{-10pt}
    }
    % Cabecera de subsección 
    \newcommand*\l@leqsubsection[2]{%
      \par\addvspace{15pt}%
      \noindent\leqIndexFont #1
      \par\vspace{-7pt}
    }
    % Título de ecuación 
    \newcommand*\l@leqt[2]{%
    \par\addvspace{10pt}%
      \par\noindent\leqTitleFont #1\par
    }
    % Miniatura de ecuación
    \newcommand*\l@leqm[2]{%
      \vspace{0.3\baselineskip}%
      \par\scriptsize\noindent\hspace*{1.5em}\(#1\)\par%
    }
    % Nota de ecuación 
    \newcommand*\l@leqnote[2]{%
      \vspace{0.2\baselineskip}%
      \par\noindent\leqNoteFont\hspace*{1.5em}\textbf{Nota.--} #1\par\vspace{0.5\baselineskip}%
    }
    % Generación del índice
    \newcommand{\mylistofequations}{%
      \clearpage
      \phantomsection
      % Título visual de la lista (ajústalo a tu gusto):
      {\centering\bfseries\Large Índice de ecuaciones\par}
      % Entrada en nuestro índice 'stc':
      \addcontentsline{stc}{section}{Índice de ecuaciones}%
      % Llamada a tu lista real (usa el nombre exacto de tu macro):
      \@starttoc{leq}%
    }
    
    %--------- Captura de títulos de sección --------%
    \let\leq@original@section\section
    \let\leq@original@subsection\subsection
    % Flag para saber si estamos en una sección asterisco
    \newif\ifLeqStarSection
    \LeqStarSectionfalse
    
    % Redefinir \section CON CONTROL DE ASTERISCO
    \renewcommand{\section}{%
      \@ifstar{\leq@section@star}{\@dblarg\leq@section@nostar}%
    }
    \def\leq@section@star#1{%
      \global\LeqStarSectiontrue
      \gdef\leq@current@section{#1}%
      \gdef\leq@current@subsection{}%
      \leq@original@section*{#1}%
      \global\LeqStarSectionfalse
    }
    \def\leq@section@nostar[#1]#2{%
      \global\LeqStarSectionfalse
      \gdef\leq@current@section{#2}%
      \gdef\leq@current@subsection{}%
      \leq@original@section[#1]{#2}%
    }
    % Redefinir \subsection
    \renewcommand{\subsection}{%
      \@ifstar{\leq@subsection@star}{\@dblarg\leq@subsection@nostar}%
    }
    \def\leq@subsection@star#1{%
      \global\LeqStarSectiontrue
      \gdef\leq@current@subsection{#1}%
      \leq@original@subsection*{#1}%
      \global\LeqStarSectionfalse
    }
    \def\leq@subsection@nostar[#1]#2{%
      \global\LeqStarSectionfalse
      \gdef\leq@current@subsection{#2}%
      \leq@original@subsection[#1]{#2}%
    }

    % ----- Restauración de valores Globales de estilo ----------%
    % Flags
    \newif\ifInFootnote
    \InFootnotefalse
    \pretocmd{\@footnotetext}{\InFootnotetrue}{}{}
    \apptocmd{\@footnotetext}{\InFootnotefalse}{}{}
    % Profundidad de listas (soporta anidadas)
    \newcount\MyListDepth \MyListDepth=0
    \AtBeginEnvironment{la}{\advance\MyListDepth by 1}
    \AfterEndEnvironment{la}{\advance\MyListDepth by -1}
    \AtBeginEnvironment{lnum}{\advance\MyListDepth by 1}
    \AfterEndEnvironment{lnum}{\advance\MyListDepth by -1}
    \AtBeginEnvironment{itemize}{\advance\MyListDepth by 1}
    \AfterEndEnvironment{itemize}{\advance\MyListDepth by -1}
    \AtBeginEnvironment{enumerate}{\advance\MyListDepth by 1}
    \AfterEndEnvironment{enumerate}{\advance\MyListDepth by -1}
    % Restauración diferida: hazlo al INICIO del próximo párrafo real
    \newif\if@pendingRestore
    \@pendingRestorefalse
    \newcommand{\RequestRestore}{\global\@pendingRestoretrue}
    \everypar{%
      \if@pendingRestore
        \global\@pendingRestorefalse
        \ifInFootnote\else
          \ifnum\MyListDepth=0
            \setlength{\parskip}{\GlobalParskip}%
            \setstretch{\GlobalBaseLineStretch}%
          \fi
        \fi
      \fi
    }
    % Al final de bloques:
    \AfterEndEnvironment{equation}{\RequestRestore}
    \AfterEndEnvironment{equation*}{\RequestRestore}
    \AfterEndEnvironment{la}{\RequestRestore}
    \AfterEndEnvironment{lnum}{\RequestRestore}
    % Al final de macros:
    \apptocmd{\eqblock}{\RequestRestore}{}{}
    \apptocmd{\imagenfit}{\RequestRestore}{}{}
    
\makeatother

% ===================== ToC propio con 4 niveles (único bloque) =====================
\makeatletter

% --- Escritores: añadimos una línea al .stc en cada cabecera numerada ---
% Guardamos las versiones actuales para llamar al comportamiento normal
\let\MyTOC@orig@section\section
\let\MyTOC@orig@subsection\subsection
\let\MyTOC@orig@subsubsection\subsubsection
\let\MyTOC@orig@paragraph\paragraph

% SECTION
\renewcommand{\section}{%
  \@ifstar{\MyTOC@section@star}{\@dblarg\MyTOC@section@nostar}%
}
\def\MyTOC@section@star#1{%
  \MyTOC@orig@section*{#1}% sin entrada al .stc
}
\def\MyTOC@section@nostar[#1]#2{%
  \MyTOC@orig@section[{#1}]{#2}% comportamiento normal (escribe al .toc)
  % Escribimos también nuestra línea al .stc (texto con \numberline)
  \addcontentsline{stc}{section}{\protect\numberline{\thesection}#2}%
}

% SUBSECTION
\renewcommand{\subsection}{%
  \@ifstar{\MyTOC@subsection@star}{\@dblarg\MyTOC@subsection@nostar}%
}
\def\MyTOC@subsection@star#1{%
  \MyTOC@orig@subsection*{#1}%
}
\def\MyTOC@subsection@nostar[#1]#2{%
  \MyTOC@orig@subsection[{#1}]{#2}%
  \addcontentsline{stc}{subsection}{\protect\numberline{\thesubsection}#2}%
}

% SUBSUBSECTION
\renewcommand{\subsubsection}{%
  \@ifstar{\MyTOC@subsubsection@star}{\@dblarg\MyTOC@subsubsection@nostar}%
}
\def\MyTOC@subsubsection@star#1{%
  \MyTOC@orig@subsubsection*{#1}%
}
\def\MyTOC@subsubsection@nostar[#1]#2{%
  \MyTOC@orig@subsubsection[{#1}]{#2}%
  \addcontentsline{stc}{subsubsection}{\protect\numberline{\thesubsubsection}#2}%
}

% PARAGRAPH
\renewcommand{\paragraph}{%
  \@ifstar{\MyTOC@paragraph@star}{\@dblarg\MyTOC@paragraph@nostar}%
}
\def\MyTOC@paragraph@star#1{%
  \MyTOC@orig@paragraph*{#1}%
}
\def\MyTOC@paragraph@nostar[#1]#2{%
  \MyTOC@orig@paragraph[{#1}]{#2}%
  % Si no numeras los \paragraph, cambia \theparagraph por texto plano sin \numberline
  \addcontentsline{stc}{paragraph}{\protect\numberline{\theparagraph}#2}%
}

% --- Impresor del índice propio (lee el .stc) ---
\newcommand{\MyTOC}{%
  \par\bigskip
  {\centering\bfseries\Large Índice de contenido\par}%
  \medskip
  \begingroup
    \parindent=0pt\parskip=2pt
    \@starttoc{stc}%
  \endgroup
}

\makeatother
% ===================== fin ToC propio =====================


%%%%%%%%%%%%%%


% --- Datos del tema ---
\title{Tema 3.A.19 -- Análisis de mercados (IV)\\
\large Diferenciación de productos: la teoría de la competencia monopolística y otros desarrollos}
\author{Marcelo Ortega}
\date{Noviembre 2025}

\oldtitle{Tema 3.A.16. Teorías del oligopolio. Análisis estático y dinámico. Teoría de juegos.}
\changerationale{La teoría de juegos como disciplina pasa a tratarse en el nuevo Tema 3.A.14. No obstante, este tema utilizará el aparato teórico de la teoría de juegos para analizar el comportamiento estratégico de las empresas en este tipo de estructura de mercado. Por lo demás, el contenido del tema no habría por qué variar sustancialmente a lo esperado de los opositores en la versión anterior del temario.}

\begin{document}


\maketitle
\changebox

\newpage

% --- Página de resumen y modificaciones ---
\puntosclave{ %% Escribir puntos clave Apuntes CECO
    }
\vspace{1cm}

\modificaciones{
    \textbf{IMPORTANTE:} La principal diferencia entre competencia monopolística y competencia en oligopolio es que en la segunda la variación conjentural NO es nula, por tanto hay siempre un cierto grado de colusión y no tienen poder de libre decisión absoluta!! Hacer mucho énfasis en esto!!

    \textcolor{red}{\textbf{OJO!}} Se puede incluir la Paradoja de COURNOT : El modelo de COURNOT asume lo que se conoce como variación conjetural nula, pero esto es inconsistente con los resultados obtenidos por el modelo, pues las curva de reacción muestran que, verdaderamente, las empresas actúan como si internalizasen los cambios del output de su contrincante ante sus propios cambios, lo que con demanda lineal sería $1/2$. Esta inconsistencia nos obliga a buscar una conjentura consistente con el comportamiento observada y, precisamente, esto nos lleva a la conclusión equivalente a la Paradoja de BERTRAND : con bienes homogeneos y competencia en cantidades, la única conjentura consistente es $\gamma_i=-1$ : solo cuando las empresas actúasen como si conjeturarán que están en competencia perfecta, sus conjenturas sobre el comportamiento del contrincante se cumplirían en un equilibrio simétrico. \href{https://www.kellogg.northwestern.edu/research/math/papers/466.pdf?utm_source=chatgpt.com}{\textit{Conjectural variations}. KAMIEN y SCHWARTZ} (Descargado)

    \textcolor{red}{\textbf{OJO!}} -- Corrección a lo anterior: Los modelos de variaciones conjenturales pueden darse por \textit{descartados}. Estas son las conclusiones del estudio anterior, si se observa con atención: las conjeturas nulas alcanzan un equilibrio equivalente al de competencia monopolística. Además, una referencia esencial para hablar de modelos conjenturales es: \href{https://www.sciencedirect.com/science/article/pii/016726819290053E}{\textit{The inconsistency of consistent conjectures: Coming back to Cournot}. LINDH (1992)}. El autor muestra todos los puntos débiles de los modelos conjenturales de forma intuitiva (aunque muy, muy compleja) y acaba concluyendo que: En aras de la mejor modelización y análisis de la competencia oligopolística, es necesario abandonar los modelos de variaciones conjenturales (y su extensión de conjeturas consistentes) con el fin de dejar paso al análisis de la interacción en un marco dinámico. No obstante, se puede preservar el concepto de conjetura (de hecho, tal y como fue propuesta por FRISCH) dentro del marco dinámico o de juego secuencial: la variación conjenturada del output rival ante variaciones nuestras. 
    }

\newpage
\MyTOC
\newpage

\mylistofequations
\clearpage

\MyListOfFigures
\clearpage


\pagenumbering{arabic}
\setcounter{page}{1}


\section{Introducción}\label{intro}


\subsection{Contextualización}\label{context}


\subsubsection{Relevancia}\label{importance}


\subsubsection{Historia}\label{history}



\subsubsection{Problemática}\label{problem} 




\subsection{Estructura de la exposición}\label{structure}






\section{Teoría del oligopolio}
\puntosclave{
    En el monopolio un producto, el equilibrio de mercado viene dado por la maximización del beneficio por parte del monopolista, siéndole indiferente para ello elegir precio o cantidad. Esta maximización se produce cuando iguala sus ingresos y costes marginales.
}
	
¿Qué es un oligopolio? Es de sobra conocido que se trata de una estructura de mercado que contiene un número reducido de empresas y un ejemplo de competencia imperfecta. Pero, ¿cómo se distingue de la competencia monopolística y/o otras estructuras de mercado en competencia imperfecta? A modo de aproximación práctica, podemos ilustrar las diferentes estructuras de mercado en función del número de competidores y el grado de diferenciación entre éstos. 

\imagenfit{EstructurasMercados.png}{Estructuras de mercado en función del grado de competidores y de diferenciación}{Elaboración propia}

\subsection{Rasgos definitorios}

Ahora bien, ¿cómo se proyecta esto sobre el desarrollo teórico-conceptual y el análisis de esta estructura por comparación al resto? Resolver esta cuestión va a ser esencial en el desarrollo de la exposición, con el fin de presentar el oligopolio y los resultados competitivos de una estructura de mercado tal.

Esencialmente, se trata de estudiar como inciden las variables que diferencias las estructuras, en el comportamiento de las firmas que operan en el mercado. 

\subsubsection{Número de competidores: Internalización de las decisiones óptimas rivales}
El primer determinante es el número de competidores: ¿qué implica que el número de competidores sea reducido? ¿cómo se proyecta esta variable en el comportamiento de las firmas en el mercado? 

El oligopolio (o competencia oligopolística, se utiliza indistintamente) se distingue del resto de estructuras de mercado porque la empresa internaliza estratégicamente la reacción de rivales identificables: toma en consideración las decisiones de sus rivales (por ejemplo, precios y cantidades) en la formación de sus decisiones óptimas.\footnote{%
    En competencia monopolística, en cambio, cada empresa ignora estratégicamente el impacto de sus decisiones sobre el comportamiento de las demás. Cada firma se comporta como un pequeño monopolista de su variedad, enfrentando una demanda residual bien definida, pero tomando como dados algunos determinantes relevantes del mercado, como el índice de precios y/o el número de empresas. Esto implica que, incluso fuera del margen, la empresa no contempla reacciones estratégicas de rivales concretos. En este sentido, la competencia monopolística es un entorno de interdependencia estratégica débil o despersonalizada: existe competencia, pero no rivalidad estratégica directa.

En DIXIT–STIGLITZ, por ejemplo, una empresa sabe que bajar su precio aumenta su demanda, pero no internaliza que ese cambio altera el índice de precios ni la demanda de las demás variedades.
} Esto es una consecuencia directa de un reducido número de competidores: si el mercado estuviera atomizado, sería imposible internalizar las decisiones óptimas de sus rivales en la formación de sus decisiones, y habría de recurrirse a mecanismos alternativos, como por ejemplo, los determinantes agregados del mercado (por ejemplo, índice de precios y número de competidores). 

En este contexto, aunque en el equilibrio la derivada marginal de la acción del rival respecto a la propia sea cero, la empresa internaliza que el entorno competitivo está formado por un número pequeño de rivales relevantes. Por eso hablamos de competencia oligopolística: el comportamiento estratégico es central y no desaparece aunque las conjeturas marginales sean nulas en equilibrio.\footnote{%
    A lo largo del tema no se va a desarrollar analíticamente el concepto de \textbf{variación conjetural}. Sin embargo, las profundas implicaciones de esta obliga a estar familiarizado con su derivación analítica: [Ver Anexo \ref{anx:variacion_conjetural}]
}

Por tanto, en oligopolio, la interdependencia estratégica es explícita y esencial: cada empresa reconoce que su decisión afecta a sus rivales y que estos reaccionan. Esto se formaliza mediante un equilibrio donde la mejor respuesta de cada empresa depende de las estrategias de las demás y es la razón por la cual las herramientas principales de análisis provienen del marco de la Teoría de Juegos. 

\subsubsection{Interdependencia estratégica: Diferenciación}
Ahora ya hemos explicado por qué la competencia oligopolística se diferencia de otras estructuras de mercado en tanto que las empresas internalizan las decisiones de sus rivales, lo que es consecuencia directa del número de competidores reducido del mercado. 

Ahora bien, imaginémos dos mercados oligopolísticos: (i) el del ecosistema de los Videojuegos, donde existen tres grandes competidores (Microsoft, Sony y Nintendo) altamente diferenciados, con ecosistemas completos con características de bien privado, servicios distintos y marcas claramente diferenciadas; y, (ii) el del mercado del petroleo donde existe un número reducido de competidores, con una diferenciación mínima (precio y cantidad disponible). 

\imagenfit{Ej_OLIGOPOLIOS.png}{Revenue total de dos mercados oligopolísticos: Ecosistemas de harwarde de videojuegos (consolas) y producción de crudo}{\textcolor{red}{Diversas}}

¿Sería correcto pensar que las dos estructuras deben estudiarse bajo la misma óptica? ¿Qué interdependencia tienen los competidores entre sí en estos mercados? ¿La misma? 

La respuesta es claramente no. ¿Por qué? Porque, aun siendo ambos oligopolios, la naturaleza de la interdependencia entre las empresas es distinta, y eso obliga a utilizar marcos analíticos diferentes. Dicho de otra forma, el grado de diferenciación entre los competidores incide directamente en la aproximación conceptual a través de la cuál analizar los mercados oligopolísticos. En términos generales, la competencia oligopolística puede estudiarse desde dos vertientes: juegos no cooperativos y juegos cooperativos.

Cuando la diferenciación es elevada, como en el ecosistema de los videojuegos, la competencia se articula en muchos márgenes (precio, catálogo, calidad, servicios, exclusividades, efectos de red, etc.). En este contexto, coordinarse resulta complejo y poco verosímil, ya que no existe una única variable competitiva sobre la que acordar. Cada empresa toma decisiones anticipando las reacciones de sus rivales, pero sin cooperación, y el resultado se describe mediante un equilibrio de Nash: la interdependencia se manifiesta como \textbf{rivalidad estratégica}.

Cuando el producto es prácticamente homogéneo, la competencia se concentra en precio y cantidad, lo que hace posibilita la coordinación entre productores. En estos mercados, el análisis relevante es el de los juegos cooperativos, donde la cuestión central no es si todos ganan con el acuerdo, sino si este puede sostenerse en el tiempo. La interdependencia es directa: la producción de cada país afecta al precio común y, por tanto, a los ingresos de todos, generando incentivos permanentes a desviarse del acuerdo.

Por tanto, el oligopolio no constituye un único modelo, sino una familia de situaciones estratégicas cuya lógica depende de dos dimensiones: (i) el grado de interdependencia efectiva entre los competidores; y, (ii) la viabilidad de sostener acuerdos sobre variables competitivas. De ahí que, sin pérdida de generalidad, el grado de diferenciación determine la forma de interdependencia y, con ello, el enfoque teórico adecuado: oligopolios diferenciados se analizan como juegos no cooperativos; oligopolios con productos homogéneos, como juegos cooperativos.



\textcolor{red}{\textbf{HASTA AQUÍ ESTÁ PERFECTO: INTRODUCIR MÁS SI ES DIDÁCTICO, SI NO ES SUFICIENTE}}





En la literatura económica actual, la competencia oligopolística se estudia a través de un conjunto relativamente bien delimitado de familias de modelos, cada una diseñada para capturar un aspecto específico de la interdependencia estratégica entre pocas empresas. No son modelos aislados, sino paradigmas que se reutilizan y adaptan en múltiples campos (organización industrial, comercio internacional, regulación, economía digital).

En primer lugar, los modelos clásicos no cooperativos: Cournot, Bertrand y Stackelberg. En segundo lugar, están los modelos de oligopolio con diferenciación de producto, que amplían Bertrand y Cournot introduciendo preferencias diferenciadas. Aquí destacan los modelos lineales tipo Hotelling y, más modernamente, los modelos de demanda discreta (logit y nested logit). Un tercer bloque clave lo constituyen los modelos dinámicos de oligopolio, donde las decisiones actuales afectan al estado futuro del mercado: modelos de inversión estratégica, I+D, aprendizaje, reputación y acumulación de capacidad, formalizados como juegos dinámicos de Markov. Finalmente, un área central en la literatura reciente es la de oligopolio con información imperfecta y contratos. Aquí se estudian problemas de colusión tácita, precios dinámicos, discriminación de precios, contratos verticales y plataformas multilaterales. 




\subsection{¿Competencia en precios o en cantidades?}
Sin embargo, la demanda residual, sobre la que ostenta poder la empresa, varía según se compita en precios (BERTRAND) o en cantidades (COURNOT).
	En la competencia à la BERTRAND. En la competencia en precios, las empresas toman como dados los precios de las demás empresas. Esto es, cuando consideran un cambio en el propio precio, cada empresa supone que las demás empresas mantendrán los mismos precios, lo que implica que las cantidades ofertadas se ajustarán al cambio en el propio comportamiento: si los bienes son sustitutivos cercanos, cada empresa anticipará que, fijando un precio ligeramente inferior al de sus rivales, podrá hacerse con todos los consumidores de sus rivales. Ello, a su vez, implicará una competencia intensa y márgenes precio-coste muy bajos. 
	En la competencia a la COURNOT. Como veremos, en la competencia en cantidades, las empresas toman como dados los outputs de las demás empresas. Ello significa que, cuando consideren distintas elecciones de cantidad a producir, cada empresa supondrá que las otras empresas ajustarán sus precios de modo que vendan las cantidades que han decidido ofertar. 

\subsubsection{Competencia y equilibrio competitivo: asignaciones eficientes}
Esta diferencia fundamental en las formas de competencia implica que la competencia conduce normalmente a precios más bajos en BERTRAND que en COURNOT.
Para entender el por qué, obsérvese que, en COURNOT los precios de los competidores deben igualar cualquier cambio en el precio de cada empresa. Por lo tanto, cuando una empresa considera un cambio en su precio, la elasticidad relevante es inferior que en el caso de BERTRAND: si la empresa i elige disminuir su precio, y espera que sus competidores mantengan su producción, conducirá a que los precios de sus competidores disminuyan, lo que a su vez atenúa el impacto del cambio en el propio precio sobre el output de la empresa i. 

Con un bien homogéneo, los resultados de un mercado en oligopolio tenderán a unos resultados más cercanos a la competencia perfecta en función de la medida en que las empresas se adapten a los precios respectivos de sus competidores. Cuando las empresas se adaptan (en primer lugar) a los precios respectivos se tenderá a asignaciones más eficientes en el sentido de PARETO, porque significa que las empresas perciben un incentivo (mayor) a reducir el precio con tal de quedarse con la demanda residual de otro competidor (undercut each other), mientras que si se adaptan en términos de cantidades ofertadas, podemos esperar que el comportamiento agresivo de una empresa conduzca a precios inferiores para todos (por lo que no habrá los misms incentivos a realizarlos).

Veamos como se refleja esto con los Modelos disponibles.


\section{Modelos no-cooperativos}
\puntosclave{
    Cuando las empresas oligopólicas compiten en un solo periodo y en cantidades (COURNOT), el precio de equilibrio se sitúa de modo realista entre el de competencia perfecta y el de monopolio.
    
	Existen dos tipos de ineficiencias asociadas al equilibrio de COURNOT: una de corte asignatario global y otra de tipo productiva, o de la industria.
    
	Cuando las empresas oligopólicas compiten en un solo periodo y en precios (BERTRAND), es preciso distinguir, en función de los costes de las empresas:
    \begin{la}
        \item Con rendimientos constantes a escala e idénticos costes marginales entre las empresas se alcanza el resultado competitivo (paradoja de BERTRAND); o,    
        \item Con rendimientos decrecientes a escala y racionamiento de la demanda, los precios de equilibrio exhibirán un comportamiento cíclico (ciclos de EDGEWORTH).
    \end{la}
	
	El modelo de KREPS-SCHEINKMAN (dinámico) permite aunar, bajo determinados supuestos, el realismo de la premisa de la competencia en precios con el realismo de los resultados de equilibrio bajo competencia en cantidades.
    
	En un marco estático, una solución cooperativa entre las empresas no es sostenible como equilibrio.
    
	El modelo de liderazgo en cantidades de STACKELBERG muestra que con dos empresas y un horizonte de dos periodos, el que una empresa pueda elegir su producción antes que la otra conduce en equilibrio a un resultado más competitivo que bajo COURNOT.
    
	Con un horizonte finito y un número de empresas fijo que tienen idénticos costes marginales constantes, la competencia en precios no logra escapar la paradoja de BERTRAND.
    
	En cambio, si el marco anterior pasa a ser de horizonte infinito o finito, pero de duración incierta, el equilibrio puede consistir en:
    \begin{la}
        \item el equilibrio bajo horizonte finito; o,
        \item estrategias condicionadas o que muestren dependencia histórica (coo las estrategias de gatillo), pero son sólo sostenibles las soluciones cooperativas entre empresas (esto es, la colusión) si el factor de descuento de las empresas es lo suficientemente alto.
    \end{la}	
}

Iniciamos el estudio del oligopolio abordando la vertiente no-cooperativa. Desde un punto de vista teórico, como Economistas, nos interesa estudiar un elemento crucial: \textbf{el equilibrio}. En particular, las propiedades del equilibrio: existiencia, unicidad y estabilidad. A partir de este análisis positivo podremos entonces comparar los resultados con la situación de competencia perfecta. 

Ahora bien, ¿cómo se alcanza este equilibrio? ¿qué es el equilibrio? El equilibrio es una condición derivada de la interacción. Sin interacción, como es evidente, no podemos hablar de equilibrio. Por tanto, la pregunta importante es la siguiente: Si para estudiar y valorar el equilibrio debemos evaluar la interacción, pero no hay, \textit{a priori}, interacción explícita en un desarrollo no-cooperativo, ¿de dónde surge la interacción y cómo podemos estudiarla?

En el análisis de mercados, la interacción surge de la competencia. Por tanto, definir cómo pueden competir las empresas será el punto de partida. Las variables estratégicas fundamentales son dos: precios y cantidades; y desde los inicios de la Teoría del Oligopolio son precisamente estas las que han centrado la mayor atención entre los economistas al estudiar la interacción en mercados oligopolísticos.

Suponiendo que: (i) existe un número finito de empresas ($I=\{1,2,\dots,I\}$); (ii) que disponen de la misma información de mercado; (iii) que enfrentan estructuras de coste similar (no necesariamente iguales); y, (iv) sin limitaciones de capacidad. Presentaremos las dos aproximaciones por orden cronológico.

\subsection{Competencia en cantidades}
La competencia en cantidades fue estudiada inicialmente por el economista y matemático francés AUGUSTIN COURNOT (1838).\footnote{%
    Se puede leer un resumen de la historia de la Teoría del Oligopolio en [\authorfont{Ver Anexo \ref{anx:historia}}]
} 

El razonamiento sería (parecido) al siguiente: si $\#I$ empresas concurrentes deciden sus niveles de producción de manera tal que maximice sus beneficios y teniendo en cuenta que todas ellas aportan sus niveles de producción al mercado; entonces emergerá un precio, de la interacción entre la oferta y la demanda, tal que se cumpla la condición de vaciamiento del mercado y el equilibrio quede garantizado.

\subsubsection{Desarrollo}
Suponiendo que cada empresa enfrenta el siguiente problema de maximización:

\eqblock{
\begin{aligned}
    &\boxed{\begin{aligned}
        &\max_{x_i}\Pi_i(x_i)=x_i\cdot P^d_x(x_i+x_{-i})-C_i(x_i)\\[0.5em]
        &\text{s.a.}\qquad \Pi_i(x_i)\ge 0
    \end{aligned}}\qquad \text{si,}\;
    \left\{\begin{aligned}
        &\Pi_i(x_i)\in\mathcal C^2\quad \text{y,}\quad \frac{\partial^2 \Pi_i(x_i)}{\partial^2 x_i}<0\\[0.5em]
        &\frac{\partial P_x^d(\cdot)}{\partial x}<0
    \end{aligned}\right.\\[2em]
    &\text{Solución:}\\
    &\over\begin{aligned}
        \\
        &\text{CPO}:\qquad \frac{\partial \Pi_i}{\partial x_i}=0\Longrightarrow x_i\cdot \frac{\partial}{\partial x_i}P_x^d(x_i+x_{-i})+P_x^d(x_i+x_{-i})-\frac{\partial}{\partial x_i}C_i(x_i)=0\\[1em]
        &\text{En Equilibrio:}\qquad IMg(x_i)=CMg(x_i)\;\Longrightarrow\;\text{Definimos:}\;\;\boxed{F(x_i,x_{-i})=IMg(x_i, x_{-i})-CMg(x_i)=0}\\[1em]
        &\overset{\text{\scriptsize Por el TFI}}{\Longrightarrow}\;\exists R_i(x_{-i})=x_i\quad \text{tal que,}\;\;\boxed{F(R(x_{-i}), x_{-i})=0}\;\;\;\forall i\in I
    \end{aligned}
\end{aligned}
}{Problema de maximización del productor: Funciones de reacción. Soluciones no-cooperativas: Modelo de COURNOT}

Por tanto, la función de reacción de la empresa $i$ obtenida mediante la derivación del problea de maximización del beneficio, internalizando las decisiones rivales y aplicando, en última instancia el Teorema de la Función Implícita, especifica la mejor acción ($x_i$) de la empresa $i$, dadas las acciones de sus competidores ($x_{-i}$). Por construcción, las funciones de reacción cumplen las condiciones de HAHN, que  garantizan que éstas exhiban una serie de propiedades deseables de cara al estudio del equilibrio: 
\begin{lnum}
    \item \textbf{Decrecientes}. La pendiente de la función de reacción es no positiva  o las cantidades son sustitutos estratégicos:  un aumento en x_i, reduce el beneficio marginal de la empresa i, lo cual conduce a la empresa i a elegir un x_i, más bajo;  y, 
    \item \textbf{Únicamente definidas}. Lo cual garantiza la continuidad de la función (la discontinuidad puede dar lugar a la inexistencia de equilibrio, tal y como muestra TIROLE, 1988).
\end{lnum}
	
Por otro lado, una característica relavante de esta formulación es que, a partir de la CPO se definir el Índice de LERNER:

\eqblock{
\begin{aligned}
    &\text{Sea:}\qquad X=x_i+x_{-i}\Longrightarrow P_x^d(x_i+x_{-i})\equiv P_x^d(X)\\[1em]
    &\Longrightarrow x_i\cdot\frac{\partial P_x^d(X)}{\partial x_i}+P_x^d(X)-C_i'(x_i)\underset{\cdot\frac{1}{P_x^d(X)}}{\Longrightarrow}\underbrace{\frac{x_i}{P_x^d(X)}\cdot \frac{\partial P_x^d(X)}{\partial x_i}}_{=\frac{1}{\varepsilon_{p,x}}}+\frac{P_x^d(X)-C_i'(x_i)}{P_x^d(X)}\\[1em]
    &\overset{\text{\scriptsize Como, $\frac{\partial X}{\partial x_i}=1$}}{\Longrightarrow} \frac{\partial P_x^d(X)}{\partial x_i}=\frac{\partial P_x^d(X)}{\partial X}\Longrightarrow \frac{x_i}{P_x^d(X)}\frac{\partial P_x^d}{\partial X}+[...]\overset{\cdot\frac{X}{X}}{\Longrightarrow} \frac{X}{X}\frac{x_i}{P_x^d(X)}\frac{\partial P_x^d}{\partial X}+\frac{X}{X}\cdot \frac{P_x^d(X)-C_i'(x_i)}{P_x^d(X)}\\[1em]
    &\Longrightarrow \frac{s_i}{\varepsilon_{p,X}}+\frac{P_x^d(X)-C_i'(x_i)}{P_x^d(X)}\underset{\varepsilon_{p,X}<0}{\Longrightarrow}\;\boxed{L_i\equiv\frac{P_x^d(X)-C_i'(x_i)}{P_x^d(X)}=\frac{s_i}{\varepsilon_{p,X}}}\\[2em]
    &\text{Reordenando $L_i$:}\qquad C_i'(x_i)=P_x^d(X)\left[1-\frac{s_i}{\varepsilon_{p,X}}\right]\overset{\text{si, }s_i>s_{j}}{\Longrightarrow}C_i'(x_i)<C_{j}(x_{j}),\quad \forall i,j\in I,\;i\neq j
\end{aligned}
}{Derivación del Índice de LERNER. Soluciones no-cooperativas: Modelo de COURNOT}

Así, el Índice de LERNER de la empresa $i$ es:
\begin{lnum}
    \item Proporcional a su cuota de mercado ($s_i$), que es endógena; y,
    \item Inversamente proporcional a la elasticidad de la demanda ($\varepsilon_{p,X}$).\footnote{%
        De hecho, la elasticidad de la oferta de la empresa $i$ puede interpretarse como la elasticidad de la demanda residual o percibida por la empresa $i$, y la demanda residual de una empresa se puede definir como la demanda a la que se enfrenta a dicha empresa, que se obtiene restando la demanda de mercado, la cantidad ofertada por los competidores a un precio dado.
    }
\end{lnum}

Además, reordenando el ínidice, vemos que las empresas gozan de poder de mercado en relación a su tamaño, lo cual a su vez está directamente relacionado con su eficiencia productiva. 

\subsubsection{Implicaciones de Política Económica}
Si suponemos que existen únicamente dos empresas (duopolio) y representando el equilibrio gráficamente, superpuesto sobre los equilibrios en estructuras de competencia perfecta y monopolio, vemos claramente las implicaciones, tanto normativas como positivas del equilibrio alcanzado:

\textbf{INTRODUCIR GRAFICO CURVAS DE REACCIÓN Y EQUILIBRIO SUPERPUESTO}

Desde un punto de vista positivo, COURNOT demuestra que el equilibrio existe, es único y estable cuando las empresas compiten en cantidades y los bienes son sustitutos. De hecho, el autor extiende dicho análisis a bienes complementarios (llegando a las mismas conclusiones) y a la competecia en precios, aqune sin resultados concluyentes. Desde una óptica normativa, la competencia oligopolística reduce la apropiación total del excedente bruto del consumidor, aumentando el bienestar en relación con la situación de monopolio. De hecho, lo anterior es conceptualmente la conclusión más importante de esta aproximación: una solución cooperativa en la que las partes se repartan la producción que maximice conjuntamente su beneficio (actuando como un único monopolista) no es sostenible como equilibrio porque cada empresa tiene incentivo para desviarse de la misma, de modo parecido a lo que ocurre en el dilema del prisionero.

\subsection{Competencia en precios: BERTRAND}
Si bien las aportaciones de COURNOT fueron el pistoletazo de salida de la Teoría del Oligopolio y, más en general, el análisis de la Competencia Imperfecta, nadie prestó especial atención a sus aportaciones hasta casi medio siglos más tade, cuando BERTRAND criticó su obra alegando que las conclusiones analíticas de COURNOT (que el equilibrio existe, es único y estable) no eran válidas.

Tras proclamar que la decisión óptima (y obvia) para los oligopolistas era la colusión, BERTRAND argumenta que las estrategias relevantes para las empresas eran los precios y no la producción.

\subsubsection{Desarrollo formal}
Suponiendo que cada empresa enfrenta el siguiente problema de maximización:

\eqblock{
\begin{aligned}
    &\boxed{\begin{aligned}
        &\max_{p_i}\Pi_i(p_,p_{-i})=p_iD_i(p_i,p_{-i})-C_i(D_i(p_i,p_{-i}))\\[0.5em]
        &\text{s.a.}\qquad \Pi_i(x_i)\ge 0
    \end{aligned}}\qquad \text{si,}\;
    \left\{\begin{aligned}
        &\Pi_i(x_i)\in\mathcal C^2\quad \text{y,}\quad \frac{\partial^2 \Pi_i(x_i)}{\partial^2 x_i}<0\\[0.5em]
        &\frac{\partial D_i(p_i,p_{-i})}{\partial p_i}<0
    \end{aligned}\right.\\[2em]
    &\text{Solución:}\\
    &\over\begin{aligned}
        \\
        &\text{CPO}:\qquad D_i(p_i,p_{-i})+(p_i-C_i(D_i(p_i,p_{-i})))\frac{\partial}{\partial p_i}D_i(p_i,p_{-i})=0\\[1em]
        &\text{Definimos:}\;\;\boxed{F(p_i,p_{-i})\equiv IMg(p_i, p_{-i})-CMg(D(p_i,p_{-i}))=0}\\[1em]
        &\overset{\text{\scriptsize Por el TFI}}{\Longrightarrow}\;\exists R_i(p_{-i})=p_i\quad \text{tal que,}\;\;\boxed{F(R(p_{-i}), p_{-i})=0}\;\;\;\forall i\in I
    \end{aligned}
\end{aligned}
}{Problema de maximización del productor: Funciones de reacción. Soluciones no-cooperativas: Modelo de COURNOT}

Por tanto, la función de reacción de la empresa $i$, obtenida mediante la derivación del problea de maximización del beneficio internalizando las decisiones rivales sobre precios (y su consecuencia sobre la propia demanda) y aplicando en última instancia el Teorema de la Función Implícita, especifica la mejor acción ($p_i$) de la empresa $i$, dadas las acciones de sus competidores ($p_{-i}$). Igual que antes, por construcción, las funciones de reacción cumplen una serie de propiedades deseables de cara al estudio del equilibrio: 
\begin{lnum}
    \item \textbf{Crecientes}. La pendiente de la función de reacción es no decreciente: un aumento en $p_i$, reduce el beneficio marginal de la empresa $i$, lo cual conduce a la empresa $i$ a elegir un $p_i$, más bajo;  y, 
    \item \textbf{Únicamente definidas}. Lo cual garantiza la continuidad de la función (la discontinuidad puede dar lugar a la inexistencia de equilibrio, tal y como muestra TIROLE, 1988).
\end{lnum}
	
También podríamos resolver, a partir del desarrollo de la CPO, el Índice de LERNER (o BERTRAND-LERNER) en este escenario. Si bien es cierto que no se utiliza, por eso no lo presentamos, es útil para ilustrar una diferencia clave con el anterior: en COURNOT la empresa internaliza cuotas de mercado, lo que se ve reflejado en la aparición de la cuota propia como componente del Índice de LERNER; mientras que en BERTRAND la empresa internaliza solo su elasticidad demanda precio directa (propia). Es una diferencia sutil, sin muchas implicaciones, pero importante a la hora de comparar.

\subsubsection{Implicaciones de Política Económica}
Si suponemos que existen únicamente dos empresas (duopolio) y representando el equilibrio gráficamente veremos claramente las implicaciones tanto normativas como positivas del equilibrio alcanzado:

\textbf{INCLUIR GRÁFICO SOLUCIÓN DE BERTRAND}

Resulta evidente que, en primer lugar, el único equilibrio estable es el de competencia perfecta. Esto es así porque, si ambas empresas tienen la misma estructura de costes, en equilibrio ambas producirán en la escala mínima eficiente, donde el coste marginal es igual o superior al coste medio y el beneficio marginal social se iguala al coste marginal social. Si, por el contrario, una empresa tuviese una estructura de costes ligeramente mejor (con costes marginales ligeramente inferiores) entonces esta empresa expulsaría a su rival del mercado y no nos encontraríamos en un dupolio. Por tanto, este equilibrio: (i) es un Equilibrio de Nash, ya que ninguna empresa tiene incentivos a desviarse; y, (ii) es equivalente al equilibrio competitivo ($P=CMg$). Esto es lo que se conoce como: Paradoja de BERTRAND. 

\subsection{Extensiones}
Ahora bien, ¿es este tipo de competencia realista? A priori puede parecer que sí, pero veámos cómo podemos salir de esta aparenre Paradoja imponiendo algunas supuestos adicionales.

\paragraph{Producción: Restricciones de capacidad}
En primer lugar, ¿es realista suponer que un competidor puede, de un instante a otro, cubrir toda la demanda de mercado que \textit{roba} a su rival mediante una rebaja en el precio? 

EDGEWORTH combatió enérgicamente la idea de que el equilibrio está determinado en condiciones de oligopolio. Para demostrarlo se basa en modelo de BERTRAND e impone limitaciones de capacidad a la producción de cada empresa (sería imponer costes marginales crecientes en su forma extrema, aunque se llega a las mismas conclusiones con costes de producción cuadráticos). Bajo esta estructura, el precio de mercado seguirá un comportamiento cíclico, experimentando subidas y caídas, comportamiento que se conoce como \textbf{ciclos de Edgeworth}.

\textbf{INCLUIR GRÁFICO CICLOS DE EDGEWORTH}

Como la colusión no constituye un equilibrio estable, ya que las empresas tienen incentivos para romper la cooperación; eventualmente, una de las empresas decidirá reducir sus precios e incrementar la producción con el objetivo de ganar cuota de mercado que hasta ahora pertenecía al competidor. La otra empresa hará lo mismo y este proceso seguirá hasta el punto en el que se alcance el máximo de la capacidad de producción, momento a partir del cual los precios se mantendrán estables ($p_{max}$). Debido a que el incremento de la demanda que sigue a una reducción de precios no vendrá satisfecha con una mayor cantidad producida, los precios comenzarán a aumentar poco a poco y las empresas verán de nuevo como sus beneficios aumentan. Con el tiempo este proceso se verá repetido y los precios oscilarán entre $p_M$ y $p_{max}$.

Otro ejemplo sería pensar en dos hoteles. Cada uno enfrenta una restricción de capacidad. En periodos de alta demanda, si un hotel sube el precio frente al precio competitivo, es posible que tenga beneficios positivos enfrentándose a una demanda residual, ya que el otro hotel no puede reaccionar (aumentando la oferta) porque ha llegado a un límite de capacidad.


\paragraph{Diferenciación}
Si el producto fuera ligeramente diferenciado (habitualmente modelizado en la literatura mediante costes de transporte adoptando un enfoque direcciones [Ver Tema 3.A.18]), entonces las curvas de reacción tendrían la siguiente forma:

\textcolor{red}{\textbf{FALTA INCLUIR DESARROLLO}}
    

\subsection{Juegos secuenciales: STACKLEBERG}
El modelo de liderazgo abordado inicialmente por HEINRICH VON STACKELBERG (1934), supone la aplicación de la competencia de COURNOT a un horizonte de dos periodos.

Este modelo se caracteriza por contar con dos empresas:
	Una empresa, llamada líder, que tiene la ventaja de elegir en primer lugar la cantidad que quiere producir;  y,
	Una segunda empresa, llamada seguidor, que no toma su decisión de producción hasta haber observado la cantidad elegida por el líder. 
De acuerdo con la Teoría de Juegos, el concepto de equilibrio relevante en un contexto dinámico con información completa es el equilibrio de NASH perfecto en subjuegos (ENPS). Para obtener el equilibrio de este modelo (que constituye un ENPS), se ha de proceder mediante retro-inducción, comenzando por lo que hace el seguidor (último subjuego). 
	Elección de x_S
La empresa seguidora resuelve el problema de optimización consistente en maximizar su beneficio dada la cantidad producida por el líder
 [La CPO de la solución es
 
A partir de la CPO se obtendría la curva de reacción del seguidor:
	Elección de x_L. 
La empresa líder maximiza su beneficio anticipando la repuesta del seguidor, i.e., eligiendo aquel punto de la curva de reacción del seguidor que maximiza su beneficio.
 
La CPO de la solución es
 
En esta CPO podemos distinguir dos partes:
	Una idéntica que en el modelo de COURNOT
 
	Otra que refleja que el líder anticipa su impacto en el seguidor (efecto estratégico), la cual conduce a producir más que en COURNOT 
 ]
En concreto, el líder sabe que al aumentar su producción, la del seguidor disminuirá (por la pendiente negativa de su curva de reacción). 
1.- El seguidor elige su output en función de su curva de reacción; 
2.- El líder, anticipando este hecho, escoge aquel output que maximiza su beneficio a lo largo de la curva de reacción del seguidor. Esto es, aquel output que resulta de la tangencia entre su línea isobeneficio más baja y la curva de reacción del seguidor.   
Por lo tanto, cabe realizar la siguiente comparación entre el equilibrio de STACKELBERG y el de COURNOT: 
	El líder produce más que bajo COURNOT, mientras que el seguidor menos (ídem con los beneficios); 
	Existencia de efecto procompetitivo. La producción total es mayor que bajo COURNOT (con lo que el precio es inferior) y, por lo tanto, la situación es mejor para el consumidor, siempre que
 
De este modo, llegamos al Equilibrio de Stackelberg, que tiene las siguientes propiedades:
	Es un equilibrio de Nash perfecto en subjuegos, ya que ninguna empresa tiene incentivos a desviarse de su estrategia.
	No es óptimo en el sentido de Pareto porque, (entre otras cosas) la seguidora podría obtener mayores beneficios a través de la cooperación sin disminuir los beneficios de la líder.
	Si los costes marginales de ambas empresas son iguales y la demanda es lineal la empresa líder produce la cantidad de monopolio (aunque los beneficios son inferiores) y la seguidora responde conforme a su función de reacción. En este caso, sí que podemos afirmar que la empresa líder produce más que la seguidora y obtiene mayores beneficios.

3.1.2.2. Juegos Bayesianos
Hasta ahora hemos supuesto que cada empresa tenía información perfecta sobre los costes de su rival (las funciones de costes son de dominio público). Levantando este supuesto damos cabida a una hipótesis más realista: la empresa no conoce la función de costes de su competidora. En un contexto con información privada (información asimétrica), ceteris paribus, resolveremos el problema suponiedo que pueden ser bajos o altos y que  la empresa no informada, inferirá las probabilidades de que se dé cada estado de la naturaleza [Resolución en Notas Finales ].
Llegaremos a un Equilibrio de Bayes Perfecto (EPB) donde la empresa sin información se verá forzada a decidir una estrategia independiente de los costes de su competidora (pero adoptará ésta siendo consciente de que el comportamiento de su competidora dependerá de sus costes). Por lo tanto, el perfil de estrategias del equilibrio será: (x_L^* (c_B ),x_L^* (c_H ),x_S^*).
En relación a los beneficios, con información incompleta ocurre lo siguiente:
	Si los costes de la empresa informada fueran bajos: c_B. 
	La empresa informada sale perdiendo con el déficit de información, es decir, le convendría más que la información fuese completa, de modo que la otra empresa conociera su punto fuerte consistente en costes bajos. 
	De manera análoga y por simetría, la empresa no informada sale ganando respecto a la situación de información perfecta.
	Los consumidores también salen ganando (disminución del precio).
	Si los costes de la empresa informada son altos: c_H, un razonamiento análogo nos mostraría que se producirían los resultados opuestos:
	La empresa informada sale ganando con el déficit de información, puesto que la otra empresa no conoce precisamente ese punto débil de la primera.
	De manera análoga y por simetría, la empresa no informada sale perdiendo respecto a la situación de información perfecta.
	Por lo que respecta a los consumidores, también salen perdiendo (aumento del precio).
 
3.4. Alternancia de modelos: adecuación e implicaciones
Como hemos visto, los modelos de BERTRAND y COURNOT no deben considerarse como modelos rivales que conducen a predicciones divergentes en cuanto al resultado de la competencia en un mercado dado, sino que se supone que describen mercados con estructuras de costes distintas: el modelo de BERTRAND puede ser una mejor aproximación a las industrias con costes marginales razonablemente constantes; el modelo de COURNOT puede ser mejor para aquéllas con costes marginales marcadamente crecientes.

La competencia en cantidad puede considerarse en general como competencia en la elección de la escala de producción, en la que la elección de la escala determina su función de coste y, por lo tanto, las condiciones de la competencia en precios. 

La idea de considerar juegos de dos fases en los que las empresas toman una decisión de inversión y, a continuación, eligen un precio no se limita a la elección de la capacidad. De hecho, en la teoría de la diferenciación de productos y en los modelos de entrada, acomodación y salida se estudian juegos de dos fases en los que la decisión de inversión atañe a una elección en el espacio de productos (p.ej., localización). Estos juegos presentan rasgos similares.  Dichos juegos en dos fases resultan atractivos porque formalizan la idea de que: las decisiones de inversión se realizan normalmente antes de las decisiones de precios.



\section{Modelos cooperativos}
\puntosclave{
    Existe una amplia literatura económica que explora los factores que inciden en la sostenibilidad de la colusión empresarial.
    
	En un contexto dinámico que permite la entrada y salida de empresas, se puede analizar el comportamiento estratégico anticompetitivo por el cual una empresa con poder de mercado trata de incidir en sus competidores de dos formas básicas:
    \begin{la}
        \item Incrementando sus costes, es decir, la creación de barreras de entrada estratégicas; o,
        \item Mermando sus ingresos, es decir, mediante la predación de precios.
    \end{la}
	
	Por lo que se refiere a las barreras de entrada estratégicas, estas normalmente se traducen en una sobre inversión en capacidad (o en dar la impresión de que esta sobre inversión se ha llevado a cabo) y solo tiene sentido cuando se precisa incurrir en un coste irrecuperable para competir.
    
	Por lo que respecta a la predicación, su éxito depende de su credibilidad, lo cual a su vez descansa en la explotación de información imperfecta en el mercado por parte de la incumbente.
}

En el apartado precedente hemos visto como nos podemos servir de la Teoría de Juegos no-cooperativos para analizar los comportamientos estratégicos de empresas que no quieren o, mejor dicho, no pueden llegar a acuerdos para restringir la competencia. En este apartado veremos las herramientas de la Teoría de Juegos cooperativos que nos sirven para analizar las situaciones en las que diferentes empresas coluden con tal de limitar la competencia de forma que el conjunto salga beneficidado. De esta forma, podemos definir la colusión como un acuerdo o práctica concertada entre empresas rivales para restringir la competencia y que se puede alcanzar de dos maneras:
\begin{la}
    \item \textbf{Explícita}. De forma explícita, mediante un acuerdo expreso entre empresas que compiten en el mismo mercado (carácter horizontal), más conocido como cártel;\footnote{%
        El ejemplo de cartel más célebre es la OPEP. Los cárteles de carácter privado suelen estar prohibidos. De hecho, el nacimiento del derecho moderno de la competencia en los Estados Unidos se debió a las protestas por las profusión de cárteles en la economía. Sin embargo, ello no impide que existan cárteles legales como: (i) los de exportación, como en Estados Unidos y Japón; y, (ii) en agricultura, por ejemplo, la PAC.
    }
    \item \textbf{Tácita}. De forma tácita, mediante una coordinación informal que se basa en el entendimiento compartido de que la cooperación redunda en el beneficio de todas las partes implicadas y que no requiere comunicación entre las partes.\footnote{%
        Desde un punto de vista jurídico, la colusión tácita no es ilegal, porque: (i) no hay acuerdo; (ii) no hay voluntad, simplemente se hace lo que se ve a nivel individual para maximizar los beneficios. Por lo tanto, se precisa evidencia material para sancionarla. En cambio, si la colusión versase sobre aplastar a los nuevos entrantes, sí sería ilegal.
    }
\end{la}

Independientemente del tipo de colusión que estudiemos, el rasgo definitorio de la cooperación es el \textit{oportunismo}: \textbf{debe ser posible cooperar para cooperar}, válgase la redundancia. Esto es clave puesto que, como decíamos al principio, la colusión solo se va a dar en mercados donde se pueda limitar la competencia a través de variables claras, identificables, cómunes y estratégicas. Por tanto, el objetivo prinicpal en esta parte no será estudiar como las empresas \textit{internalizan} las decisiones estratégicas de sus rivales sino cómo se configura (o debe configurarse) la cooperación de forma que se garantice su estabilidad o, al menos, se minimicen los incentivos a desviarse. 



\subsection{Colusión explícita: Alianzas sostenibles}
La primera gama de herramientas que estudiaremos pretenden establecer las condiciones para que una alianza sea sostenible. De esta forma, la colusión tiene su origen en un pacto o acuerdo explícito sobre una o varias variables estratégicas, un reparto de las ganancias totales del mercado y la minimización de los incentivos a desviarse o formar subgrupos alternativos.

Para amenizar la presentación podemos recurrir a un ejemplo práctico sobre el cártel más célebre del mundo: la OPEC. Supongamos que somos uno de los principales exportadores de petróleo en la década de 1960: demanda mundial creciente y elevada dependencia energética de las economías industrializadas. Un poder de negociación débil frente a las grandes compañías petroleras y los países importadores conduce a precios bajos y rentas disipadas. En este escenario, cualquier intento de acuerdo explícito enfrenta problemas clásicos de colusión: incentivos individuales a sobreproducir, riesgo de rupturas internas, heterogeneidad en costes y necesidades fiscales dispares. ¿Existe alguna herramienta para analizar los beneficios de una alianza y distribuirlos de manera tal que se minimicen los riesgos de desviación?Cómo nos ayuda? Presentaremos dos enfoques alternativos.

\subsubsection{Valor del núcleo}
El primero es el valor del núcleo, que se define como el conjunto de asignaciones estables que un subconjunto de agentes no puede mejorar. Por tanto, un requisito lógico para que un vector de asignaciones sea mutuamente aceptable es que ningún subgrupo pueda mejorar rompiendo el acuerdo. Formalmente:
  
Se trata de el conjunto de resultados donde ningún Jugador tiene incentivos para desviarse, es decir, dejar la coalición y formar su propia coalición. El concepto de estabilidad es muy parecido al concepto de Equilibrio de Nash. 
EJEMPLO
 	Imaginemos un Juego en el que tres jugadores pueden unir sus fuerzas para conseguir 100€. Dos jugadores juntos pueden alcanzar un pago comprendido de 20€ y un jugador en solitario puede alcanzar 0€. Veamos el desarrollo análitico de este juego (izq.)
Sin embargo, cabe señalar dos problemas en cálculo del valor del núcleo:
	En algunos casos, el núcleo puede ser el conjunto vacío, de forma que no existe una asignación que sea estable ante desviaciones de subgrupos;
Sólo en escasas ocasiones el núcleo contendrá un único vector.

\subsection{Valor de SHAPLEY}
Por otro lado, tenemos el valor de SHAPLEY.

El Valor de SHAPLEY es un método de distribución de riquezas que se basa en el reparto esperado de cada agente en base a un mismo Juego cooperativo (o no) bajo todo el \textit{scope} de configuraciones posibles. El razonamiento es como sigue: una coalición de jugadores coopera, y obtiene una cierta ganancia general de la cooperación. Dado que algunos jugadores pueden contribuir más a la coalición que otros, o pueden poseer diferente poder de negociación (por ejemplo, amenazando con destruir todo el excedente), ¿qué reparto final de los beneficios de la cooperación entre los jugadores debemos esperar que surjan en cualquier juego en particular? O expresado de otra manera: ¿qué importancia tiene cada jugador para la cooperación global, y qué recompensa puede él o ella razonablemente esperar? 

El valor de Shapley ofrece una posible respuesta a esta pregunta. Formalmente:
 
Por lo tanto, el Valor de Shapley nos proporciona una manera muy comprensiva de calcular el resultado justo que debería obtener el Jugador i dentro de una gran coalición (es decir, la coalición formada por todos, esto es, la coalición N), como función de sus contribuciones marginales tanto a la coalición como al Juego en su conjunto. La fórmula se puede interpretar de la siguiente manera: imaginar que la coalición está formando un actor a la vez, con cada actor exigiendo su contribución v (S ∪ {i}) - v (S) como una compensación justa, y luego para cada actor tomar el promedio de esta contribución sobre las diferentes posibles permutaciones en la que se puede formar la coalición. 

4.1.5.1. Propiedades
El Valor de Shapley tiene las siguientes propiedades deseables:
	Eficiencia. La ganancia total se distribuye de manera que el sumatorio de los valores de Shapley para cada miembro de la gran coalición es igual a la ganancia de la gran coalición. 
 
	Simetría. Si dos jugadores juegan papeles idénticos en la gran coalición, i.e. son equivalente, entonces sus valores de Shapley también serán los mismos. 
Si  , entonces  
	Linealidad o aditividad. Para un mismo conjunto de jugadores N, el juego consecutivo de cualesquiera dos juegos v y w, y el juego compuesto v + w, reportan el mismo pago conjunto para cada coalición. Esto es, la distribución de pagos debe ser invariante respecto a una descomposición arbitraria del Juego. 
  
También por cada número real a:  
	Zero-player o jugador ficticio. El valor de Shapley que contribuya a la coalición exactamente con su pago individual se le debe asignar este pago. También se puede describir como, el valor de Shapley de un jugador nulo (aquel que recibe de forma esperada un pago cero) debe tener esta asignación en el resultado final.
Habiendo presentado estos dos conceptos, la diferenciación de la Teoría de Juegos cooperativos no significa que la Teoría no cooperativa no sea capaz de explicar la cooperación entre grupos de individuos. De hecho, ahora pasaremos a estudiar como esta otra rama de la Teoría de Juegos nos muestra cómo la cooperación puede surgir fruto del comportamiento racional, en ausencia de capacidad para realizar acuerdos vinculantes que es un supuesto de partida en la Teoría de Juegos cooperativos.

\subsection{Colusión tácita}
Pasamos ahora cómo podemos las diferentes formas de garantizar el equilibrio colusorio cuando este emerge de la interacción entre los agentes (tácitamente). 

La manera más intuitiva de reconocer este equilibrio es pensar como sería la interacción entre agentes dentro de un juego que son plenamente conscientes se va a repetir. De primeras, es normal pensar que este conocimiento conduce de manera inequívoca a la colusión como mínimo en la primera repetición del juego, por ser ésta la única forma de apropiarse de obtener un beneficio extraordinario (por comparación a la no colusión: competencia perfecta). Veámos si el desarrollo analítico acompaña nuestra intuición. 

4.1.1. Supuestos
El escenario de horizonte infinito se estudiará bajo los siguientes supuestos simplificadores:  
	La competencia es en precios, existen dos empresas (los resultados serían fácilmente extensibles a n empresas),
	Las empresas tienen el mismo coste marginal constante c,
	El juego no varía de un periodo a otro, dado que no se producen ni inversiones ni cambios en el entorno competitivo de la empresa. Este supuesto permite considerar el juego total como la suma de juegos independientes (e iguales) en cada periodo. 
A diferencia del caso anterior, la empresa no sabe nunca si se encuentra en el último periodo, por lo que no se puede proceder mediante retroinducción para obtener el ENPS.

4.1.3. Desarrollo: FRIEDMAN
A continuación, pasamos a mostrar por qué los precios cooperativos o colusorios constituyen un ENPS, siguiendo el análisis pionero de FRIEDMAN (1971).
4.1.3.1. Tipos de estrategias
En este caso, existen dos ENPS del superjuego:

De estrategias incondicionada. El Jugador opta por la decisión “no-cooperativa” (Pase lo que pase, opto por no cooperar). Este tipo de estrategia no muestra dependencia histórica y consiste en la repetición del Equilibrio de Nash (estático), al igual que bajo un horizonte finito. 
 Estrategias condicionadas. Consistentes en un perfil de estrategias cuyo equilibrio son precios colusorios. Son un tipo de estrategias que muestran dependencia histórica, ya que especifican el precio en el periodo t_n como función de la historia -periodos t_0,t_1,…,t_(n-1). Se trata del mantenimiento de cualquier precio tal que:
 
Donde p^cp es el precio bajo el equilibrio estático  y p^M el precio bajo monopolio. Por ello, denominaremos estos precios p de equilibrio como precios cooperativos o colusorios, dado que se basan en que todas las empresas del mercado cooperan fijando un precio supracompetitivo (i.e., superior a aquél bajo el equilibrio estático). Existen diferentes tipos de estrategias a este respecto, entre las que podemos distinguir:
	Trigger-strategies. (Original en el trabajo de JAMES FRIEDMAN, 1971) Se inicia el Juego con un precio cooperativo (consciente o inconscientemente) y se mantiene mientras ninguna de los integrantes modifique su precio. Si este cambio se da, se desencadena una espiral de bajada como herramienta de castigo (en caso de que una empresa no haya seguido el precio cooperativo o colusorio en el pasado). Se conocen en la literatura como trigger strategies (estrategias de gatillo) porque una sola desviación desencadena el castigo y un paro en la cooperación entre las empresas a la hora de fijar precios supracompetitivos.  Por lo tanto, como la opción dos es un castigo creíble (ninguna empresa puede salir ganando desviándose unilateralmente del mismo), es preciso comprobar cuando la opción uno constituirá un equilibrio. En último instancia se llega al Equilibrio de Nash no cooperativo p^cp.
	Tit-for-tat. Esta estrategia se basa en la idea de que un jugador debe cooperar en la primera ronda y luego hacer lo mismo que hizo el otro jugador en la ronda anterior. Es decir, si un jugador coopera en la primera ronda, el otro jugador también cooperará. Si un jugador no coopera en una ronda, el otro jugador tampoco cooperará. 
Éste autor considera que la siguiente estrategia condicionada de cualquier empresa es consistente cuando fija los siguientes precios de la siguiente forma:


4.1.3.4. Mantenimiento del Equilibrio cooperativo: valor descuento
El mantenimiento del equilibrio cooperativo o colusorio será sostenible si el beneficio intertemporal de ceñirse al precio cooperativo es superior al beneficio intertemporal de desviarse del mismo:
 ≥ 
donde los periodos van de 0 a infinito, δ es el factor de descuento (igual y constante para todas las empresas) ,  y Π^C es el beneficio de la empresa bajo cooperación, Π^R es el beneficio que obtiene una empresa al desviarse del precio cooperativo o colusorio y Π^P es su beneficio bajo el castigo.  Dado que toda desviación iniciará una guerra de precios con vuelta al equilibrio estático, la mejor desviación será aquella que maximice el beneficio a corto plazo; es decir, fijar un precio ligeramente por debajo del cooperativo o colusorio y así acaparar toda la demanda. Ahora bien, también requiere que la amenaza de llevar a cabo la estrategia trigger sea realmente creíble. 

Entonces, podremos mostrar que la cooperación es sostenible si el factor de descuento es lo suficiente alto (δ es superior a un cierto valor). Es decir, si las empresas son lo suficientemente pacientes (perciben que el largo plazo es importante):


    
4.2.3.5. Gráficamente
Se puede apreciar de la siguiente forma, de modo que el valor de los beneficios futuros descontados de la colusión sea mayor que el de la desviación, por lo que el área roja debe ser menor que el área azul (los incentivos a desviarse no son suficientes):
 
 
Cabe destacar que el resultado de FRIEDMAN  entronca con la contribución pionera de STIGLER (1961), quien destacó cómo la búsqueda de mayores beneficios detrás de la colusión conduce también a que las empresas quieran traicionar al cártel y acaparar el mercado (de modo análogo a lo que cabe observar en el dilema del prisionero). Esto ha dado pie a una amplia literatura económica que explora los factores que hacen más sostenible la colusión, ya sea por incidir en: (i) el beneficio colusorio; (ii) el beneficio que resulta de desviarse; o, (iii) el beneficio bajo el castigo, que no solo depende del tamaño del castigo infligido, sino también de la rapidez con que se detecta el incumplimiento del pacto colusorio por parte de una empresa. (Una información más detallada de cada tipo de variable que condiciona la sostenibilidad de la colusión se puede ver en [Ver Anexo I])
Aunque pueda dar la impresión de que la colusión sólo sería sostenible en un horizonte infinito. Sin embargo, la colusión también sería posible en un horizonte finito pero de duración incierta. Ello se puede modelizar suponiendo que en cada periodo existe una probabilidad x de que las empresas sigan compitiendo en el siguiente periodo.  Se puede mostrar que esta situación es equivalente a un horizonte infinito, pero con un factor de descuento inferior (x-δ). 

4.2.4. Críticas
De todas formas, volviendo a nuestro razonamiento inicial, cabe señalar que el análisis de la interacción repetida ha sido objeto de crítica precisamente porque no especifica cómo se ponen de acuerdo las empresas a la hora de escoger el precio colusorio y la cuota de mercado que atiende cada empresa (como se ha visto, existen múltiples equilibrios colusorios).  ¿No resultaba intuitivo? Precisamente una contracrítica podría ser decir que aquí subyace una demostración de que el agente \textcolor{red}{\textbf{sí parece ser racional, tal y como postula la Teoría Económica ortodoxa}}.

En la literatura se supone normalmente que se escoge el precio más alto y que el mercado se reparte a partes iguales (pero ello no resulta válido en el caso de asimetrías en costes), hacer cumplir el resultado colusorio y, en concreto, elegir el castigo en caso de incumplimiento por parte de una empresa. 

\subsection{Implicaciones de Política Económica}
Desde un punto de vista positivo, hemos visto que el equilibrio existe, es único y estable en cualquiera de las dos vertientes. 



Desde un punto de vista normativo, el perjuicio que causan las conductas anticompetitivas a los consumidores en una economía de mercado (principalmente los cárteles, pero también las barreras estratégicas erigidas por empresas con elevado poder de mercado para dificultar artificialmente la entrada en su entrada y así proteger sus beneficios supracompetitivos, [Ver Apartado 5]) ha conducido a que casi todos los países castiguen y persigan este tipo de prácticas empresariales. 
En España, las sanciones económicas por conductas anticompetitivas son las más altas de nuestro ordenamiento jurídico, lo cual se queda aún corto frente a las penas de cárcel con las que se llega a sancionar a los directivos que participen en un cartel en los países anglosajones.

\subsubsection{Factores que fomentan la cooperación}
La colusión puede afectar muchas dimensiones del comportamiento de las empresas: la segmentación de mercados, la cantidad ofertada, las inversiones a realizar,\footnote{%
    Eligiendo las inversiones de tal modo que suavice la competencia, normalmente limitando las inversiones
} pero la variable que se suele elegir es: el precio.

La literatura económica ha señalado diferentes características de una industria o factores que afectan a la sostenibilidad de la colusión, a través de su impacto en el factor de descuento necesario para mantener la colusión (afectando a los beneficios de deserción o la magnitud y probabilidad del castigo): 
	El número de empresas. Cuanto mayor es el número de empresas presentes en el mercado, más dificil se hace mantener la colusión (se gana lo mismo que antes desertando, pero no se gana tanto bajo colusión porque el beneficio se reparte entre más empresas); 
	La transparencia del mercado. Que afecta la habilidad para detectar defecciones al inicidir sobre: (i) el tiempo de detección; y, (ii) la cantidad de información asimétrica , 
	La frecuencia de interacción en el mercado. Cuando las empresas compiten entre sí con frecuencia es más fácil sostener la colusión porque las empresas pueden reaccionar más rápido ante una desviación de una de ellas (adelantar el castigo). ,  
	El crecimiento de la demanda. Facilita la colusión, ya que el coste de la deserción es mayor; 
	Las asimetrías en el coste. Redundan en una mayor dificultad en mantener la colusión porque: (i) un desertor eficiente puede ganar más con la deserción; y, (ii) resulta más dificil castigar a un desertor eficiente.  
	La presencia simultánea en varios mercados (multimarket contact). La presencia en distintos mercados facilita la sostenibilidad de la colusión por diferentes motivos: (i) incrementa la frecuencia de interacción entre las empresas; (ii) puede permitir a las empresas sostener la colusión en mercados en los que las características de la industria por sí solas no lo permitirían, castigando la desviación en todos los mercados;  (iii) puede permitir suavizar las asimetrías que surjan en mercados individuales.  La presencia en ambos mercados restablece una simetría general que facilita la colusión.


\clearpage
\section{Comportamientos \textit{ad-extra}: Barreras de entrada}

\textcolor{red}{\textbf{ESTE APARTADO, SI SE QUIERE INCLUIR, DEBE SER MUUUUUY BREVE}. Correponde a otro tema y no es necesario incluirlo aquí}


Hasta aquí hemos estudiado los comportamientos \textit{ad intra} de las estructuras de mercado oligopolítisticas. Ahora bien, siempre y cuando no existen barreras de entrada, los incumbentes deberán también tener presente en todo momento la posibilidad de que otras empresas decidan entrar en el mercado y, en consecuencia, mermar su poder de mercado. La Teoría del Oligopolio también ha dado origen a diversos análisis sobre prácticas comunes cuyo objetivo es disuadir a los posibles entrantes y son las que presentaremos brevemente en este apartado. 

En general, distinguimos dos formas básicas de disuasión:
\begin{la}
    \item \textbf{Barreras de entrada estratégicas}. Las barreras de entrada estratégicas consisten en inversiones realizadas por el incumbente, en una fase previa a la competencia, que persiguen incrementar artificialmente los costes de sus competidores, ya sea actuales o potenciales. Por lo tanto, en este caso, la inversión actúa como variable estratégica. Normalmente, esta barrera se traduce en una sobreinversión en capacidad: reducción de costes, calidad, publicidad , etc.\footnote{%
        Evidentemente, solo tiene sentido hablar de las barreras estratégicas cuando se incurre en un coste irrecuperable (sunk costs), puesto que de no ser así, la escala mínima eficiente \textit{estaría} en cualquier lugar de la recta de Costes Marginales (lo óptimo sería, evidentemente el mínimo, pero se alcanzaría \textit{obligatoriamente}).
    };  y,
    \item \textbf{Predación}. Persigue disminuir de modo artificial los ingresos del entrante, con el fin de que decida que no es rentable continuar en el mercado.
\end{la}

\subsection{Barreras de entrada estretégicas: Precio límite}
El modelo más famoso de barreras de entrada estratégicas es el de precios límite (BAIN, 1956; SYLOS-LABINI, 1962; MODIGLIANI, 1958), cuya idea principal es que, bajo ciertas condiciones, la empresa incumbente puede sostener un precio tan bajo que disuade la entrada en el mercado. Su justificación es que, al observar un precio reducido, el entrante potencial esperará márgenes reducidos en el mercado y, por lo tanto, beneficios más bajos. Este argumento fue objeto de controversia  hasta que SPENCE (1977) y DIXIT (1979, 1980), por un lado, y MILGROM y ROBERTS (1982), por otro, clarificaron sus rasgos subyacentes:


En términos generales, la reformulación debida a SPENCE y DIXIT considera el modelo de STACKELBERG de competencia secuencial en cantidades como uno de elección secuencial de capacidades. Esto es, aunque la competencia en el mercado (si la hay) determine el precio de mercado a corto plazo, en el largo plazo las empresas compiten a través de la acumulación de capital (ya hemos visto antes, al esbozar el modelo de KREPS-SCHEIKMAN, que las cantidades se pueden reinterpretar como capacidades).
 
La ventaja de estar ya presente en el mercado (la posibilidad de acumular capital de antemano) conduce a la empresa incumbente a acumular una gran capacidad (y, por lo tanto, a fijar un precio reducido) de modo a disuadir o limitar la entrada. ,  
Sin embargo, este modelo puede ser criticado porque la amenaza de la incumbente no sea creíble: si la nueva empresa entra, no sería racional que la incumbente mantuviera la producción inicial, ya que lo lógico sería acomodar la entrada de la nueva empresa adaptando su producción para maximizar beneficios.

5.1.1.1. Precio límite: asimetrías informativas
La reformulación del argumento del precio límite debida a MlLGROM y ROBERTS se basa en la existencia de asímetrías informativas entre la empresa instalada (incumbente) y la entrante. El incumbente fija un precio bajo, no porque disponga de una elevada capacidad de producción, sino porque trata de transmitir la información de que, ya sea la demanda o su coste marginal son bajos, señalizando con ello al potencial entrante una baja rentabilidad si entra. 

\subsection{La predación}
Por el contrario, la predación es una figura mucho más compleja, bajo la cual se incluye todo aquella acción agresiva y \textbf{costosa} por parte de una empresa que persigue mantener o aumentar su poder de mercado.\footnote{%
    Las acciones son: (i) agresivas, porque dañan directamente a los competidores obligándoles a abandonar el mercado, reducir su inversión, escala de actividad o abanico de productos; y, (ii) costosas, porque se trata de una acción subóptima hoy (corto plazo) con el propósito de aumentar el poder de mercado en el futuro y obtener beneficios que compensen las pérdidas en que se pueda incurrir a corto plazo.
}
	
La predación adopta primordialmente la forma de precios reducidos.  Por lo tanto, la propia naturaleza de esta práctica resulta desconcertante: en principio, los precios reducidos se asocian normalmente con un mayor bienestar social, pero, en ciertas ocasiones, tienen un fin anticompetitivo. En estos casos:
	A corto plazo. Los precios reducidos mejoran el bienestar a corto plazo, durante el tiempo que dura la predación; pero,
	A largo plazo. Una vez la presa (entrante) abandona el mercado, el precio aumenta, por lo que el efecto final (si tiene éxito) será perjudicial para el bienestar a largo plazo, al eliminar la competencia. 
 
5.2.1. Condicionantes del éxito
El factor que determina que una empresa racional decida salir del mercado son los beneficios futuros. Por lo tanto, el éxito de la predación depende de su credibilidad: el hacer creer que los escasos beneficios actuales se prolongarán el día de mañana.
Por tanto, el elemento clave que permite la credibilidad de la predación es la existencia de información imperfecta (incertidumbre), ya que, caso contrario, la predación sería observable y no creíble (el entrante, o la entidad financiera que financie al entrante, se percataría de que el incumbente no puede llevar a cabo un comportamiento subóptimo indefinidamente).  
5.2.2. Tipos de predación
Ello se evidencia en los dos tipos básicos de predación que se distinguen en la literatura: 
	La predación financiera. Persigue deteriorar la situación financiera de la presa (por ejemplo, reducir su cash-flow) para impedir que obtenga financiación para acometer sus inversiones, ya que la entidad financiera utilizará el cash-flow generado como indicio de la rentabilidad del proyecto (riesgo percibido del préstamo). El origen de este tipo de predación se encuentra en el argumento conocido como the long-purse story (o deep-pocket story).  Esta teoría se ha reelaborado modernamente para dotarla de bases teóricas más sólidas, fundamentándola en la existencia de imperfecciones en el mercado de crédito (debido a asimetrías informativas por parte de la entidad financiera);  ya que, de otro modo, la predación sería observable y la empresa podría obtener la financiación 
	Manipulación de la información. Sobretodo acerca de la rentabilidad del mercado mediante el precio (señal),  confundiendo al competidor acerca de: 
	La dureza del incumbente (información sobre uno mismo). Dando la impresión de agresivo e irracional en el mercado, de modo que no sólo se pretende expulsar al primer entrante, sino también disuadir a cualquiera que se lo plantee; 
	La demanda del mercado. Dando la impresión de que la demanda es tan baja que sólo hay sitio en el mercado para una empresa, y así retrasar su entrada;  
	La eficiencia en costes del incumbente. 


El modelo de Precio Límite debido a MILGROM y ROBERTS [Ver Apartado 5.1.1.1.] también se puede reinterpretar en términos de predación mediante manipulación de la información (sobre costes o demanda del mercado, vid. supra), si se persigue inducir a la salida del mercado. ,  


\subsection{Implicaciones de Política Económica}
El perjuicio que suponen estas prácticas empresariales (unilaterales) para el bienestar ha motivado que también sean sancionadas por las autoridades de competencia (en la UE bajo la figura del abuso de posición dominante y en Estados Unidos, bajo la figura de la monopolización del mercado). 

Cabe destacar también que estos modelos que se van a presentar, relativos a las disuasión que pueda ejercer una/s empresa/s incumbente frente a posibles entrantes, han permitido rechazar la visión simplista de la Escuela de Chicago [Ver Tema 3.A.15] (basada en información perfecta) según la cual los recortes de precios son siempre respuestas naturales a shocks de costes o de demanda, o una mayor presión competitiva.









\clearpage
\section{Conclusión} 

\subsection{Extensiones y relación con otras partes del Temario}


\subsection{Opinión}



\subsubsection{Idea final (Salida o cierra)}


\clearpage
\pagenumbering{gobble}
\MyBibliography
\MyBibItem{Autor}{Título}{Año}{DOI -- LINK}




%ANEXOS
\clearpage
\anexo{\textbf{Preguntas Test}}
%\input{\TemaCode/questions.tex} 


\clearpage
\anexo{Historia: Teoría del Oligopolio}\label{anx:historia} 

\href{https://antonibosch.com/uploads/book_file/498_EC-VIVES_Capitulo1.pdf}{\textit{Precios y Oligopolios. Ideas Clçasicas y Herramientas Modernas. Capítulo I -- El origen de la Teoría del Oligopolio y de la Teoría moderna de los Juegos}. X. VIVES, 2001}

Cournot, Bertrand y Edgeworth establecieron los fundamentos de la moderna teoría del oligopolio y estudiaron el tema central de cómo se forman los precios en un mercado con un número reducido de competidores. Así precedieron a la moderna teoría de los juegos no cooperativos en la propuesta de conceptos de solución para situaciones de interacción estratégica. 

Un tema común entre los precursores de la teoría del oligopolio es la cuestión de si los precios están determinados cuando hay interacción estratégica. De hecho, plantearon la mayoría de temas que más tarde se demostrarían centrales en la teoría del oligopolio: el concepto de solución adecuado, el análisis del equilibrio (existencia, unicidad, y estabilidad), así como las propiedades de la estática y dinámica comparativa. También plantearon preocupaciones, aunque parecen haber pasado desapercibidas en la literatura, sobre la plausibilidad del concepto mismo de equilibrio. Más tarde, Chamberlin, Hotelling y Robinson enfatizaron la importancia de la diferenciación de producto. Repasemos pues brevemente sus contribuciones uno a uno.

Cournot (1838) fundó la teoría del oligopolio proponiendo un concepto de solución para la interacción oligopolística, examinando tanto el caso de bienes sustitutivos como complementarios y estudiando la estabilidad de la solución propuesta. También estudió la posibilidad de colusión y la conexión entre oligopolio y competencia perfecta. Cournot argumentó que los precios también se podían determinar en condiciones de rivalidad oligopolística. Imaginó empresas (las famosas productoras de agua mineral) compitiendo independientemente y decidiendo sus niveles de producción. Cada empresa aporta esos niveles de producción al mercado, donde surge un precio de la interacción entre oferta y demanda. El equilibrio se alcanza cuando el volumen de producción de cada empresa representa la respuesta óptima ante los volúmenes de producción de las demás empresas. Es decir, cuando cada empresa decide un nivel de producción que maximiza sus beneficios dados los volúmenes de producción de las empresas rivales. En terminología moderna esto es un equilibrio de Nash (Nash (1950)) con los volúmenes de producción como variables estratégicas. Como es bien conocido, el equilibrio conlleva un precio superior al coste marginal de producción. Nuestro autor también estudió la competencia entre empresas productoras de bienes complementarios. Es interesante señalar que en este caso Cournot supuso que las empresas competirían en precios y aplicó el mismo concepto abstracto de solución (de Nash). De nuevo, un equilibrio se caracteriza por acciones, precios en este caso, que son consistentes en el sentido que la acción de una empresa representa la respuesta óptima a las acciones de las empresas rivales. Esto es un equilibrio de Nash con precios como variables estratégicas. Además, demostró que el precio de equilibrio que surge bajo competencia en productos complementarios, contrariamente al caso de bienes sustitutivos, es mayor que el precio de monopolio o de cártel Cournot también apuntó a los fundamentos de la teoría moderna de la competencia perfecta cuando consideró la ausencia de límites, es decir, la competencia perfecta. Para que las acciones de las empresas, en el límite, no tengan efectos sobre el precio del mercado, Cournot argumentó que tenían que ser pequeñas con respecto al mercado, cada una con un volumen de producción inapreciable. 

La contribución de Cournot pasó desapercibida para los economistas hasta la reseña de su libro por Bertrand en 1883. Bertrand tras proclamar que la decisión obvia para los oligopolistas es la colusión, argumentó que en el caso de mercados de producto homogéneo considerados por Cournot las estrategias relevantes para las empresas eran los precios y no los volúmenes de producción. En lenguaje moderno, Bertrand propuso el concepto de solución de Nash en precios. Si adoptamos esta sugerencia, y los costes unitarios de producción son constantes e iguales para todas las empresas, como en el ejemplo de los productores de agua mineral de Cournot, el precio se iguala al coste marginal, la solución competitiva. 

Edgeworth combatió enérgicamente las ideas de Cournot, en particular la afirmación de que el equilibrio está determinado en condiciones de oligopolio. En su artículo sobre la teoría pura del monopolio (1925, pp.117-118, publicado originalmente en 1897) dice: 

\begin{cita}[\authorfont{EDGEWORTH}][1925]
    \textit{He }(COURNOT)\textit{ concludes that a determinate proposition of equilibrium defined by certain quantities of the articles will be reached. Cournot's conclusion has been shown to be erroneous by Bertrand for the case in which there is no cost of production; by Professor Marshall for the case in which the cost follows the law of increasing returns; and by the present writer for the case in which the cost follows the law of diminishing returns}
\end{cita}

La idea central de Edgeworth es que en situaciones de números pequeños, en oligopolio, y en contraste con el monopolio o la competencia perfecta, el equilibrio es indeterminado. Para ilustrar su teoría presenta un modelo de competencia en precios en el que un duopolio produce dos bienes sustitutivos o complementarios x e y. A partir de una función de utilidad para el consumidor lineal en el dinero, $U(x, y)$, Edgeworth define los bienes como rivales si $\frac{\partial^2 U}{\partial x\partial y}<0$ para todos los consumidores (o, por lo menos, como media), y complementarios si la derivada cruzada es positiva. A partir de la especificación de la función de utilidad, se derivan las funciones directas e inversas de demanda y se analizan sus propiedades (notando, por ejemplo, que la concavidad de la función de utilidad y bienes rivales implica sustituibilidad bruta). 

Cuando los bienes son sustitutivos Edgeworth consideró un modelo de producto homogéneo producido por un duopolio con costes marginales crecientes, bien en la forma extrema de restricciones de capacidad (1925), o con costes de producción cuadráticos (1922). En ambos casos concluyó que los precios nunca alcanzarían una situación de equilibrio, sino que oscilarían de forma cíclica indefinidamente. Edgeworth también señaló que el nivel de indeterminación disminuye conforme aumenta el grado de diferenciación de los bienes, en el límite las empresas devienen monopolios independientes. Una conclusión fundamental de su análisis es que una empresa nunca tendrá incentivos para producir por encima del nivel competitivo (determinado por su coste marginal) para cualquier precio. De acuerdo con Edgeworth, el análisis debe tener en cuenta esta restricción voluntaria al intercambio. En el caso de restricciones de capacidad, la empresa simplemente no puede ofrecer más que su capacidad disponible. La consecuencia es que el equilibrio competitivo, que es el único posible con estrategias puras en el duopolio de Edgeworth (ver Shubik (1959)), puede desestabilizarse: una empresa aumentando su precio puede aumentar sus beneficios, dado que su rival puede no ser capaz de (o no querer) abastecer todo el mercado al precio competitivo. Este hecho junto con el argumento tradicional de recorte de precios de Bertrand conlleva la no existencia de equilibrio (en estrategias puras) para valores plausibles de los límites de capacidad de las empresas o de las especificaciones de los costes. Este modelo ha dado lugar a lo que ahora se denomina competencia Bertrand-Edgeworth, donde las empresas compiten en precios, pero también donde no se exige a ninguna empresa que sirva toda la demanda al precio prevalente en el mercado. En este tipo de competencia la existencia de equilibrio, de acuerdo con el análisis moderno de la teoría de juegos, sólo está garantizada en estrategias mixtas, donde las empresas deciden de forma aleatoria sobre un cierto intervalo de precios. 

En el escenario de bienes complementarios, Edgeworth consideró como al menos muy probable que el equilibrio económico es indeterminado, a diferencia de la situación de productos rivales, donde pensó que con certeza ése es el caso. Como ilustración considera bienes complementarios perfectos producidos por empresas competidoras (en precios) y también los peculiares competidores \textit{Nansen y Johansen tirando de sus trineos }(tras haber muerto los perros)\textit{ sobre las llanuras árticas}. La razón de la indeterminación puede ser, como en el caso de los bienes rivales, que la forma de la función de beneficios es tal que un equilibrio (de Nash) no existe. Sin embargo, incluso cuando existe, el equilibrio puede no surgir, aunque ambas empresas sean perfectamente inteligentes y conscientes de prever los motivos de los rivales. En realidad, si sólo una de las empresas es perfectamente inteligente y capaz de prever las decisiones futuras, mientras que la otra empresa es miope, el equilibrio propuesto por Edgeworth es el que, más tarde, propuso Stackelberg (1934) en su modelo de liderazgo. El punto (de Stackelberg), en palabras de Edgeworth, no será un equilibrio estable, excepto bajo el supuesto extremo que Nansen sea perfectamente inteligente y capaz de prever las decisiones futuras mientras que Johansen, como dice el refrán, \textit{no ve más allá de sus narices}.\footnote{%
    En el contexto del juego de precios con bienes complementarios: la empresa capaz de prever decisiones futuras optimiza teniendo en cuenta la función de reacción de la empresa rival que es decreciente. La competencia en precios con bienes complementarios puede entenderse como la situación dual de la competencia en cantidades con productos sustitutivos, la cual normalmente presenta funciones de reacción decrecientes.
} Si ambas empresas prevén futuras acciones, entonces pueden intentar ser cada una más lista que la otra y el candidato a equilibrio (de Nash) puede desestabilizarse ya sea por accidente o porque un jugador ha previsto algún movimiento futuro fuera del equilibrio. 

Edgeworth parece, pues, apuntar a una crítica moderna del equilibrio de Nash: el supuesto de que la racionalidad (maximización de la utilidad dadas unas creencias) y los pagos son de conocimiento común no es suficiente para explicar el comportamiento de equilibrio (de Nash). La literatura sobre racionabilidad (Bemheim (1984) y Pearce (1984)) ha señalado que el equilibrio de Nash no es el único resultado que se deriva del conocimiento común de pagos y racionalidad. Las estrategias de equilibrio son racionalizables, pero en general el conjunto racionalizable es mucho mayor. Por ejemplo, en un oligopolio de Cournot con demanda lineal, con más de dos empresas y costes marginales constantes, cualquier precio entre el competitivo y el monopolístico es racionalizable. 

Resumiendo, Edgeworth pensó que el problema del oligopolio era fundamentalmente indeterminado y que los precios nunca alcanzarían posiciones de equilibrio en un mercado caracterizado por números pequeños, a diferencia de lo que ocurre en mercados competitivos. En el caso de bienes sustitutivos, su línea de ataque es señalar la inexistencia de equilibrio (en estrategias puras). Con bienes complementarios, el ataque es más fundamental dado que cuestiona el concepto mismo de equilibrio (de Nash). 

Un desarrollo reciente que arranca a partir de un enfoque diferente de la teoría del oligopolio, la teoría de los juegos supermodulares, proporciona condiciones para limitar el grado de indeterminación en el comportamiento estratégico, una preocupación central de Edgeworth. Como hemos mencionado antes, Edgeworth fue el primero en considerar una función de utilidad genérica $U(x,y)$ y en introducir supuestos sobre las segundas derivadas cruzadas. De hecho, el supuesto $\frac{\partial^2 U}{\partial x\partial y}>0$ se define, a veces, como la complementariedad de Edgeworth. En terminología moderna, cuando el pago marginal de la acción $x$ de un agente es creciente en la acción $y$ (que puede ser una acción del mismo agente o de otro agente), decimos que las acciones $x$ e $y$ son complementarias estratégicas. La teoría de los juegos supermodulares proporciona el marco teórico para estudiar las complementariedades estratégicas y está ligado al tema edgeworthiano de la indeterminación. Así, los juegos supermodulares siempre tienen equilibrios en estrategias puras, el conjunto de equilibrios tiene un orden y propiedades de estabilidad atractivas y un comportamiento racionalizable dentro de ciertos límites. 

Hotelling (1929), Chamberlin (1933) y Robinson (1933) introdujeron la diferenciación de producto en el estudio de la competencia imperfecta. Todos ellos pensaron en términos de empresas compitiendo en precios. El modelo del segmento de Hotelling incorpora consumidores con gustos heterogéneos y proporciona los fundamentos de la teoría de la localización. Chamberlin (1933) y Robinson (1933) avanzaron el modelo de competencia monopolística. En el modelo del grupo grande de Chamberlin, las empresas son pequeñas con respecto al mercado, pero retienen algún poder de monopolio. Esta es la base del modelo de competencia monopolística en el que, aunque cada empresa se enfrenta a una curva de demanda decreciente en su variedad, no puede influenciar las magnitudes agregadas del mercado. 

Dos temas más han reclamado la atención en la teoría del oligopolio. La estabilidad del equilibrio y la tendencia de los oligopolistas a colusionarse. Cournot pensó que su propuesta de solución sería estable en el sentido que una desviación del equilibrio se autocorregiría a través de una serie de reacciones de las empresas. La dinámica del desequilibrio estaría gobernada por las funciones de respuesta óptima de las empresas. Esto es, por las funciones que generan el volumen de producción óptimo de una empresa en función de los volúmenes de producción de los rivales. Este concepto de estabilidad ha sido objeto de crítica porque las empresas de Cournot en el proceso de ajuste hacia el equilibrio ven como sus expectativas son sistemáticamente rechazadas por la evidencia. Las empresas esperan que sus rivales mantengan sus niveles de producción, mientras que en realidad los varían constantemente. Sólo en el punto de equilibrio las expectativas se ven confirmadas. Esto condujo a Fellner (1949) a afirmar que Cournot \textit{tenía razón por las razones equivocadas}. Sin embargo, esta crítica no es aplicable al concepto de equilibrio. En un equilibrio de Cournot, las empresas anticipan correctamente los niveles de producción de los rivales y optimizan adecuadamente. Además, el proceso de ajuste de Cournot se ve ahora, desde una perspectiva evolucionada, como una forma rudimentaria de modelar el comportamiento de agentes con racionalidad limitada. 

La posibilidad de colusión entre productores ha sido una fuente permanente de preocupación en la teoría del oligopolio. Cournot evocó la cuestión de por qué las empresas no se ponen de acuerdo en compartir el mercado al precio de monopolio. El pensó que, desde el punto de vista de una empresa si los rivales respetan el acuerdo, esta empresa tiene interés en violarlo y, así, mejorar su nivel de beneficios. En otras palabras, tiene un incentivo para hacer trampa. Cournot pensó que el tipo de acuerdo de monopolio o cártel sólo podían mantenerse por medio de un compromiso formal (1838, p. 83). \textbf{Chamberlin} (1929) \textbf{defendió con fuerza la tendencia a la colusión en oligopolio}. Pensaba que mientras \textbf{en el caso del grupo grande, el modelo de competencia monopolística era adecuado, en el caso del grupo pequeño, las empresas se darían cuenta de su interdependencia y actuarían (implícitamente) de forma que maximizasen los beneficios conjuntos, teniendo en cuenta el posible uso de estrategias de penalización contra aquellas empresas que no respetasen el acuerdo}. Esta visión de la colusión tácita ha sido apoyada por muchos economistas influyentes, incluyendo a Fellner (1949), Samuelson (1967) y Stigler (1964) y se ha popularizado en los libros de texto de economía industrial como Scherer y Ross (1990). Stigler (1964) ya avanzó que el mayor obstáculo a la estabilidad de los acuerdos colusivos es el recorte secreto de precios. Cuando una empresa observa un bajo nivel de ventas, no está segura de si se debe a un shock adverso de demanda para su producto o a que sus rivales están violando el acuerdo de cártel. 

El \textbf{problema del oligopolio} gira alrededor de la potencial indeterminación de los precios cuando hay pocos competidores. La opinión de Cournot de que bajo condiciones de oligopolio los precios estarían determinados fue enormemente combatida por Edgeworth, quien pensaba que los precios eran esencialmente indeterminados. La opinión de Bertrand se acerca más hacia la de Cournot, mientras que Chamberlin tiende hacia Edgeworth. El concepto de solución adecuado y sus propiedades (existencia, unicidad, estabilidad y estática comparativa) permanecen como preocupaciones del análisis moderno utilizando las herramientas de la teoría de los juegos. La obra de von Neumann y Morgenstem \textit{Theory of Games and Economic Behavior} se publicó en 1944 y los artículos de Nash sobre juegos no cooperativos se publicaron en 1950 y 1951 . Aunque su influencia fue reconocida por algunos investigadores en el campo de la teoría del oligopolio (principalmente Shubik(1959)), la aplicación masiva de la teoría de juegos en el análisis de la competencia entre empresas tuvo que esperar hasta finales de los años setenta. 

Precisamente, es la interacción entre la teoría de juegos y la teoría del oligopolio lo que ha hecho posible la formalización de importantes ideas sobre competencia en contextos de organización industria. Hemos descrito las contribuciones de la teoría del oligopolio a la teoría de juegos estáticos no cooperativos, incluyendo los más recientes desarrollos de los juegos supermodulares. La teoría de juegos ha permitido desarrollar importantes herramientas para el análisis de las ideas de compromiso, interacción dinámica y situaciones de información incompleta. Para los juegos dinámicos con información completa, el concepto central es el de equilibrio perfecto en los subjuegos (EPS), un refinamiento del equilibrio de Nash debido a Selten (1965). Para los juegos de información incompleta, Harsanyi (1967-68) introdujo el concepto de equilibrio bayesiano de Nash (EBN). 

Para los juegos dinámicos de información incompleta, el concepto de equilibrio bayesiano perfecto (EBP) es un refinamiento del equilibrio de Nash que combina las ideas del EPS y del EBN con las reglas de la actualización bayesiana. Refinamientos más fuertes del equilibrio de Nash han sido propuestos por Selten (1975), con el equilibrio perfecto, y por Kreps y Wilson ( 1 982a), con el equilibrio secuencial. 

El modelo de líder-seguidor propuesto por Stackelberg (1934) incorpora la idea de compromiso al decidir el líder en primer lugar su volumen de producción anticipando la reacción del seguidor. El equilibrio de Stackelberg es un ejemplo de equilibrio perfecto en los subjuegos de Selten, en el que se eliminan las amenazas no creíbles. En el modelo de Stackelberg, por ejemplo, el seguidor no puede amenazar con inundar el mercado si el líder no produce una cierta cantidad. Esto sería una amenaza no creíble porque el seguidor no tiene incentivos para cumplir su amenaza si el líder no actúa como tal. 

Stackelberg (1952, pp. 194-195) pensó que su solución de líder-seguidor, que él denominó duopolio asimétrico, \textit{es inestable porque el vendedor pasivo puede reiniciar la lucha en cualquier momento }[...]\textit{ Es posible, por supuesto, que los duopolistas intenten suplantarse en el mercado, de manera que se inicie una competencia 'a muerte'}. Esto es lo que se denomina el \textit{punto de guerra de Stackelberg}, en el que cada empresa en un duopolio intenta ser el líder en un juego de competencia en cantidades. Pensó también que el duopolio era un régimen inestable y el equilibrio sólo se recuperaría con \textit{un monopolio colectivo o una regulación estatal}. 

Schelling (1960) fue el pionero del análisis de situaciones en las que una debilidad inicial, como la restricción sobre las estrategias disponibles para un jugador, puede convertirse en una ventaja estratégica. En un marco más general, la restricción a estrategias de Markov, es decir, estrategias que dependen únicamente de variables relevantes en términos de pagos, han demostrado ser útiles para comprender la dinámica del compromiso cuando variables de estado, como los límites de capacidad, por ejemplo, condicionan la competencia a corto plazo. Las ventajas de una empresa instalada en un mercado sujeto a entrada proporciona una ilustración clarificadora de estos fenómenos. El equilibrio perfecto de Markov (Maskin y Tirole (1997)) proporciona el concepto de equilibrio adecuado. Utilizando las herramientas que proporcionan los juegos diferenciales podemos extender el análisis al estudio de la competencia en tiempo continuo. El resultado es una rica teoría que explica una variedad de esquemas dinámicos de precios y en los que la complementariedad o sustituibilidad estratégica intertemporal juega un papel crucial. 

La teoría de los juegos repetidos, iniciada por Friedman (1971), Aumann y Shapley (1976), y Rubinstein (1979), proporciona una herramienta adecuada para analizar los mecanismos que permiten mantener acuerdos de colusión utilizando amenazas creíbles. Una extensión de la teoría, iniciada por el modelo de oligopolio de Green y Porter (1984), en la que se introduce un control imperfecto de las acciones de las empresas puede explicar el papel de las guerras de precios como argumento que permite sostener la colusión y evitar recortes secretos de precios. 

Los desarrollos producidos por el análisis de juegos de información incompleta han contribuido, de forma decisiva, al estudio de los precios cuando las empresas tienen información privada y a explicar fenómenos como la puesta en común de información y la revelación estratégica de información. El diseño de mecanismos, una aplicación de la teoría de juegos con información incompleta, ilustra las restricciones de incentivos a las que se enfrenta una empresa cuando trata de revelar información, por ejemplo, con propósitos colusivos. Los juegos dinámicos de información incompleta permiten el estudio de la inversión en una variable intangible como la información o la distorsión de la información. La utilización de los precios como forma de manipular información con sus distorsiones asociadas de señalización ayuda en la construcción de teorías sobre precios límite y el problema del polizón en problemas de barreras a la entrada, inercia en cuotas de mercado, precios para introducirse en un mercado, y comportamientos tendentes a perjudicar a competidores potenciales. 

Resumiendo, la teoría del oligopolio ha dado un impulso a la teoría de juegos no cooperativos a partir del concepto de equilibrio básico, el equilibrio de Cournot-Nash, así como su caracterización en términos de las propiedades de existencia y estabilidad. A su vez, la teoría de juegos ha proporcionado valiosas herramientas, desarrolladas en ocasiones en colaboración con aplicaciones del oligopolio, como los juegos supermodulares y los juegos repetidos con control imperfecto, que han demostrado ser fundamentales en análisis dinámicos y de situaciones de información incompleta.


\clearpage
\anexo{Variación conjentural: ¿Qué es? Derivación analítica}\label{anx:variacion_conjetural}

La derivación analítica de la variación conjentural se puede presentar de muchas maneras. En este anexo se presentarán dos formas: el origen de la variación conjentural en la competencia en precios y su origen en los modeos de competencia en cantidades.

$$\text{Variación conjentural: \textbf{PRECIOS}}$$

Denotamos $(x_1, x_2)$ la producción de las empresas 1 y 2, respectivamente. Supongamos que:

$$\begin{aligned}
    &r_1(x_1,x_2)\equiv \text{Ingresos de la empresa 1}\qquad \wedge\qquad c_1(x_1)\equiv\text{Costes de la empresa 1}\\[1em]
    &r_2(x_1,x_2)\equiv \text{Ingresos de la empresa 2}\qquad \wedge\qquad c_2(x_2)\equiv\text{Costes de la empresa 2}\\[2em]
    &\text{Teniendo en cuenta que,}\quad 
    \left\{\begin{aligned}
        \frac{\partial r_1(x_1,x_2)}{\partial x_2}\leq 0\\[1em]
        \frac{\partial r_2(x_1,x_2)}{\partial x_1}\leq 0
    \end{aligned}\right.
\end{aligned}$$

Por tanto, sabemos que hay cierto grado de sustituibilidad entre los bienes. Partiendo de estos supuestos, podemos expresar los beneficios empresariales como: 

$$\begin{aligned}
    \Pi_1(x_1,x_2)=r_1(x_1,x_2)-c_1(x_1)\qquad \wedge\qquad\Pi_2(x_1,x_2)=r_2(x_1,x_2)-c_2(x_2)
\end{aligned}$$

A esta misma ecuación se le podrían incluir impuestos (introduciendo $1-\tau$ como multiplicador de ambas \textit{curvas} de ingresos). Por último:

$$\begin{aligned}
    \gamma_1=\frac{\partial x_2}{\partial x_1}\equiv \text{Conjetura de la empresa 1 sobre la variación de la empresa 2}\\[1em]
    \gamma_2=\frac{\partial x_1}{\partial x_2}\equiv \text{Conjetura de la empresa 2 sobre la variación de la empresa 1}\\[1em]
\end{aligned}$$

Mediante la resolución de los programas de maximización del beneficio de ambas empresas, obtendríamos las siguientes condiciones de primer orden (CPO's): 

$$\begin{aligned}
    &\text{Empresa 1}:&&\frac{\partial \Pi_1}{\partial x_1}=\frac{\partial \Pi_1}{\partial x_1}+\frac{\partial\Pi_1}{\partial x_2}\frac{\partial x_2}{\partial x_1}-c_1'(x_1)\quad \Longrightarrow\quad \frac{\partial r_1}{\partial x_1}+\gamma_1\frac{\partial r_1}{\partial x_2}-c_1'(x_1)=0\\[1em]
    &\text{Empresa 2}:&&\frac{\partial \Pi_2}{\partial x_2}=\frac{\partial \Pi_2}{\partial x_2}+\frac{\partial\Pi_2}{\partial x_21}\frac{\partial x_1}{\partial x_2}-c_2'(x_2)\quad \Longrightarrow\quad \frac{\partial r_2}{\partial x_2}+\gamma_2\frac{\partial r_2}{\partial x_1}-c_2'(x_2)=0
\end{aligned}$$

Supnogamos que se cumplen las condiciones de segundo orden (suficientes) para que los resultados obtenidos sean efectivamente un máximo.\footnote{%
    Nótese que ambas CPO's representan una forma general de las funciones de reacción en la que no hemos hecho ningún supuesto, ni sobre los ingresos, ni sobre los costes, ni tan solo sobre las variaciones conjeturales.
} En este contexto (si son efectivamente un máximo), por el Teorema de la Función Implícita podemos obtener las Curvas de reacción para cada empresa dadas las condiciones óptimas derivadas de la CPO. Definimos: 

$$\begin{aligned}
    &\text{Empresa 1}:&&F_1(x_1,x_2,\gamma_1)\equiv \frac{\partial r_1}{\partial x_1}+\gamma_1\frac{\partial r_1}{\partial x_2}-c_1'(x_1)=0\\[1em]
    &&&\underset{\substack{\text{\scriptsize podemos resolver}\\\text{\scriptsize para $x_1$ como función de $x_2$}}}{\Longrightarrow}\boxed{x_1=R_1(x_2)}\\[2em]
    &\text{Empresa 2}:&&F_2(x_1,x_2,\gamma_2)\equiv \frac{\partial r_2}{\partial x_2}+\gamma_2\frac{\partial r_2}{\partial x_1}-c_2'(x_2)=0\\[1em]
    &&&\underset{\substack{\text{\scriptsize podemos resolver}\\\text{\scriptsize para $x_2$ como función de $x_1$}}}{\Longrightarrow}\boxed{x_2=R_2(x_1)}
\end{aligned}$$

En este escenario, podemos entonces conluir que ambas constituyen un sistema de ecuaciones cuya resolución nos daría el Equilibrio de Nash: $(x_1^*,x_2^*)$. Ahora bien, podemos ganar en detalle si vemos lo que veraderamente ocurre en base a sus conjenturas ya que las reacción de las empresas estará determinada por sus conjeturas:

$$\begin{aligned}
    R_1(x_2)=\frac{\partial x_1}{\partial x_2}=-\frac{\frac{\partial F_1}{\partial x_2}}{\frac{\partial F_1}{\partial x_1}}\sim\gamma_2\quad\wedge\quad R_1(x_2)=\frac{\partial x_1}{\partial x_2}=-\frac{\frac{\partial F_1}{\partial x_2}}{\frac{\partial F_1}{\partial x_1}}\sim\gamma_1
\end{aligned}$$

Por tanto, aquí vemos la clave de la variación conjentural: ¿qué reacción lleva a cabo la empresa ante una variación de la competidor? Analíticamente:

$$\begin{aligned}
    &\underset{\substack{\\[1em]\\\text{\scriptsize (equivalente para la empresa 2)}}}{\begin{aligned}
        &\text{Si}:&& \gamma_2>R_1(x_1^*)\Longrightarrow\text{1 sobreestima la reacción de 2}\\[1em]
        &\text{Si}:&& \gamma_2=R_1(x_1^*)\Longrightarrow\text{1 estima correctamente la reacción de 2}\\[1em]
        &\text{Si}:&& \gamma_2<R_1(x_1^*)\Longrightarrow\text{1 subestima la reacción de 2}\\[2em]
    \end{aligned}}\\[1em]
    &\Longrightarrow\quad \underbrace{e_1\equiv \gamma_1-R_2(x_1^*)\qquad \wedge\qquad e_2\equiv\gamma_2-R_1(x_2^*)}_{\text{\scriptsize Errores conjeturales}}
\end{aligned}$$

Observamos por tanto cómo deben interpretarse las variaciones conjenturales (respuestas en base a creencias formadas por cada competidor) y como estas dan origen a lo que se conoce como \textbf{errores conjeturales} cuyas implicaciones serán muy importantes.

Ahora bien, vemos que la incorporación de estas conjenturas es muy sútil y, explícitamente, no se puede apreciar \textit{ningún cambio}. Dicho de otra manera, el cambio experimentado se verificará comparándolo con la conjentura \textit{correcta}. Entonces, ¿podemos generalizar este marco? La respuesta es: depende. Para verificar el impacto de una conjetura \textit{correcta} debemos definir todos los fundamentos del modelo: costes, ingresos, impuestos y subvenciones si lo hay. Solo de esta forma podremos verificar el error conjetural. 

A modo de ejemplo, si suponemos una situación de duopolio en competencia à la COURNOT, con empresas simétricas, podemos representar el impacto de las conjeturas de la siguiente forma:

$$\begin{aligned}
    p=a-b(x_1+x_2),\quad c_1(x_1)=c_2(x_2)\\[1em]
    \left\{\begin{aligned}
        R_1(x_2)=\frac{a-c_1}{2b}-\frac{(1+\gamma_1)}{2}x_2\\[0.5em]
        R_2(x_1)=\frac{a-c_2}{2b}-\frac{(1+\gamma_2)}{2}x_1
    \end{aligned}\right.\Longrightarrow
    \quad \underset{\substack{\text{\scriptsize (la empresa es más aversa al riesgo: más cauta)}\\\text{\scriptsize y viceversa (si es $\gamma_j=-1$: paradoja de BERTRAND)}}}{\uparrow\gamma_j\Rightarrow\,\uparrow \frac{\partial R_j }{\partial x_{-j}}}
\end{aligned}$$

Lo primero que te diría, después de leer los textos que adjuntaste, es que tu intuición va en una dirección muy importante: el “error conjetural” no debe interpretarse automáticamente como algo perjudicial para la empresa que lo comete. De hecho, una parte sustancial de la literatura muestra precisamente que una conjetura “errónea” puede llevar a un resultado más rentable para la firma que la conjetura “correcta”; y eso no es una anomalía accidental, sino una señal de que el criterio de “corrección” de la conjetura no coincide necesariamente con el criterio económico relevante, que suele ser rentabilidad, estabilidad o poder de mercado. En otras palabras, una conjetura puede ser falsa como descripción de la reacción rival y, sin embargo, útil estratégicamente o al menos no dañina en equilibrio. Ésa es una de las razones por las que el debate sobre “conjeturas consistentes” acabó siendo mucho más problemático de lo que parecía al principio. 

Conviene separar con mucho cuidado tres cosas que en esta literatura se mezclan con facilidad. Una cosa es la conjetura en sí: qué cree una empresa que hará la rival si ella cambia su cantidad o su precio. Otra cosa es el equilibrio que resulta de que ambas optimicen dadas esas creencias. Y una tercera, mucho más exigente, es la “consistencia” de la conjetura: que la reacción que la empresa atribuye a la rival coincida con la reacción que efectivamente tendría lugar. El atractivo inicial del concepto de consistencia es evidente: Fellner criticaba que, en modelos tipo Cournot, las empresas estaban “en lo cierto por las razones equivocadas”, porque en equilibrio aciertan la producción rival, pero no aciertan cómo reaccionaría el rival si alguien se desviara. Esa crítica empujó a muchos autores a exigir que las conjeturas fueran correctas no sólo en el punto de equilibrio, sino también respecto de las reacciones fuera de él. Ése es el origen intelectual de las “conjeturas consistentes”. 

Ahora bien, cuando uno entra en serio en la literatura, aparece la primera gran conclusión: no hay un único concepto de consistencia. Lindh distingue entre una consistencia meramente puntual, que sólo exige que la conjetura sea correcta en el punto de equilibrio, y una consistencia local, que exige que lo sea en un entorno del equilibrio. La puntual es demasiado débil: apenas mejora el problema de Fellner, porque fuera del punto de equilibrio las empresas siguen “equivocándose” sobre las reacciones. La local es mucho más exigente, pero justamente por eso introduce paradojas lógicas y supuestos informativos muy fuertes. Lindh es tajante en esto: los intentos de imponer consistencia intertemporal y, al mismo tiempo, restringir los equilibrios posibles, han sido en general poco exitosos; o bien no restringen mucho, o bien descansan en una historia dinámica que resulta conceptualmente inestable.  

Éste es, a mi juicio, el núcleo del problema que tú percibes. Si una empresa supiera realmente cómo reacciona la rival en un entorno del equilibrio, entonces ya no tendría sentido modelar sus decisiones como si la cantidad rival fuera un dato exógeno a observar y luego responder. En el extremo, si la reacción rival es conocida, la empresa incorpora ese conocimiento a su propia optimización; pero la rival hace lo mismo; y entonces aparece una circularidad: cada una elige sabiendo que la otra elige sabiendo que ella elige. Lindh lo formula de manera muy fuerte: la “inconsistencia” no está realmente en el modelo estático, sino en la historia dinámica que se cuenta para justificarlo. En equilibrio, dice, las producciones esperadas se confirman mutuamente en cualquier modelo de variaciones conjeturales; el problema surge cuando se pretende dar una interpretación dinámica a esas supuestas reacciones y se exige que esa historia sea plenamente coherente. Ahí aparece la paradoja.

Dicho de forma didáctica: imagina un duopolio Cournot. Cada empresa elige cantidad. Si cada una cree que la otra no reaccionará, tienes la conjetura de Cournot, es decir, variación conjetural cero. Fellner diría: eso es ingenuo, porque si una empresa se aparta del equilibrio, la otra sí reaccionará. Hasta aquí, todo parece claro. Pero cuando intentas “corregir” esa ingenuidad y exiges que la empresa conozca la reacción verdadera de la rival, entonces la empresa deja de enfrentarse a una cantidad rival dada y pasa a enfrentarse a una regla estratégica de la rival. En ese momento ya no estás en el Cournot ordinario; estás metiendo un nivel adicional de anticipación estratégica que cambia por completo el sentido del problema. Por eso Lindh sostiene algo que parece paradójico pero es muy revelador: bajo ciertas formulaciones de consistencia local, el propio Cournot puede aparecer como “consistente”. No porque la rival no reaccione realmente, sino porque, si la trayectoria óptima ya está internalizada, la mejor respuesta local puede equivaler a no reaccionar a “tests” de la rival. El resultado suena extraño, pero el mensaje profundo es que el concepto de consistencia está semánticamente resbaladizo.

Por eso creo que la mejor respuesta a tu pregunta es ésta: sí, el análisis de errores conjeturales tiene consistencia intelectual, pero no porque permita identificar sin ambigüedad qué conjeturas son “correctas” y cuáles “incorrectas”. La tiene más bien como análisis de cómo distintas creencias estratégicas alteran el equilibrio y la distribución de beneficios. En ese sentido, la literatura de Kamien y Schwartz es muy útil. Ellos muestran que, para una estructura de mercado dada, el equilibrio puede desplazarse desde resultados cercanos a competencia hasta resultados cercanos a monopolio según cuáles sean las variaciones conjeturales. Es decir, las conjeturas no son un mero adorno psicológico: funcionan como un parámetro de conducta que modifica producción, precio y beneficios. Y cuanto más alta es la respuesta esperada del conjunto de rivales, menor tiende a ser la producción de equilibrio y mayores los beneficios. Esto encaja exactamente con tu observación: una “mala” conjetura, entendida como no coincidente con la reacción verdadera, puede empujar a la firma a una conducta más dura o más blanda que, dadas las respuestas de los demás, la sitúe en una posición mejor. 

De hecho, Kamien y Schwartz subrayan que un aumento de la variación conjetural reduce la cantidad de equilibrio, eleva el precio y aumenta los beneficios, siendo el máximo beneficio el de coordinación perfecta. Eso ya basta para ver por qué no puedes identificar “error” con “perjuicio”. Una empresa puede equivocarse sobre la mecánica exacta de reacción de sus rivales y, sin embargo, actuar de un modo que sostenga precios más altos o cantidades más bajas, obteniendo una renta superior a la que obtendría si tuviera una conjetura “correcta” en el sentido técnico de consistencia. La verdad descriptiva de la conjetura y su conveniencia estratégica no son lo mismo. 

Aquí aparece una segunda idea importante: el criterio de conjetura consistente no selecciona, en general, el resultado “económicamente razonable” en el sentido de mayor beneficio privado. Lindh insiste justamente en que, en muchos modelos de cantidad, el equilibrio con conjeturas localmente consistentes da beneficios inferiores al Cournot ordinario, y por eso no parece plausible como resultado de un proceso genuino de aprendizaje o de racionalidad estratégica. Si una empresa, aprendiendo y usando toda la información disponible, acaba en un equilibrio que le da menos beneficio que otro equilibrio accesible con creencias menos “consistentes”, cuesta mucho vender la consistencia como principio normativo o positivo de conducta. Ahí es donde tu objeción es especialmente fuerte: si la conjetura “equivocada” mejora la posición de la firma, la apelación a la consistencia pierde mucha de su fuerza económica. 

La aportación de Kamien y Schwartz matiza esto, pero no lo destruye. Ellos no dicen que la consistencia sea inútil; dicen, más bien, que bajo ciertas hipótesis sí se pueden caracterizar conjeturas consistentes. Por ejemplo, muestran que si se exige consistencia con conjeturas constantes en equilibrio simétrico, entonces la demanda debe pertenecer a una familia funcional muy concreta. En el caso de duopolio, llegan incluso a resultados llamativos: con demanda de elasticidad unitaria, cualquier conjetura inferior a uno puede ser consistente en equilibrio; y en otros casos, la consistencia de conjeturas constantes en cantidades o en precios depende finamente de la forma de la demanda. El mensaje no es que la consistencia sea un principio universal y robusto, sino que, para que exista y sea manejable, hacen falta estructuras funcionales muy específicas. Eso es exactamente lo contrario de una justificación general del concepto.  

Esto también ayuda a responder tu inquietud sobre el análisis gráfico. Si por “análisis gráfico” entiendes el clásico diagrama de curvas de reacción en el plano de cantidades, la literatura no dice que sea completamente inútil. Lo que sí dice, de manera bastante clara, es que puede ser profundamente engañoso si se interpreta de forma fuerte como representación literal de expectativas de reacción coherentes y dinámicamente fundadas. Lindh lo dice casi sin rodeos: la simplicidad intuitiva del diagrama de reacción —la idea de que basta dibujar dos curvas y ver cómo convergen— oculta una trampa lógica. La imagen de pizarra es útil pedagógicamente, pero no debería tomarse como una representación fiel de un proceso real de aprendizaje mutuo bajo información perfecta. En ese sentido, sí: si el modelo conjetural se interpreta gráficamente como si mostrara una historia dinámica real y lógicamente consistente de reacciones y contrarreacciones, la crítica de Lindh implica que ese análisis gráfico no tiene un fundamento sólido. Pero si se usa de forma modesta, como dispositivo estático para representar un sistema de mejores respuestas percibidas, sigue siendo útil




































\end{document}